\documentclass[11pt]{article}
\usepackage[paperwidth=210mm,paperheight=297mm,margin=16mm]{geometry}
\usepackage{fontspec,amsmath,array,multirow,graphicx}
\usepackage{polyglossia}
\setmainlanguage{latin}\setotherlanguages{arabic,greek}
\setmainfont{Linux Libertine O}[Ligatures=TeX]
\newfontfamily\arabicfont{Amiri}[Script=Arabic]
\newfontfamily\greekfont{FreeSerif}[Script=Greek]
\newfontfamily\NallinoSigns{NallinoSigns.otf}[Path=./fonts/]
\newcommand{\AbjadZero}{{\NallinoSigns\symbol{"E000}}}
\newcommand{\AbjadOver}[1]{\leavevmode\vbox{\hrule height .45pt\kern1.2pt\hbox{#1}}}
\newcommand{\Name}[1]{{\addfontfeatures{LetterSpace=12}#1}}
\pagestyle{empty}\setlength{\parindent}{0pt}
\renewcommand{\arraystretch}{1.18}

\begin{document}
\clearpage
{\large \textarabic{٢٤٥}}\par\vspace{2mm}
\begin{center}\textbf{Fol. 226,r.—237,r.}\end{center}
\begin{Arabic}\begin{tabular}{|p{84mm}|>{\centering\arraybackslash}p{9mm}|>{\centering\arraybackslash}p{9mm}|>{\centering\arraybackslash}p{9mm}|>{\centering\arraybackslash}p{9mm}|>{\centering\arraybackslash}p{14mm}|>{\centering\arraybackslash}p{17mm}|}\hline
\multicolumn{7}{|c|}{\parbox{168mm}{\centering\large ابتداء جداول اسماء الكواكب الثابتة ومواضعها لسنة \AbjadOver{اقضا}\textsuperscript{1} لذي القَرْنَين}} \\ \hline
\multirow{2}{*}{\parbox{80mm}{\centering من اسماء الكواكب الثابتة التي في الصُّوَر الشماليّة من مِنْطَقَة البروج}} & \multicolumn{2}{c|}{الطول} & \multicolumn{2}{c|}{العرض} & \multirow{2}{*}{\parbox{13mm}{\centering\small علامات الجهة}} & \multirow{2}{*}{\parbox{16mm}{\centering\small مراتب العظمة}} \\ \cline{2-5}
 & {\small درج} & {\small دقائق} & {\small درج} & {\small دقائق} & & \\ \hline
\multicolumn{7}{|c|}{\large من الدُّبّ الأصْغَر\textsuperscript{2}} \\ \hline
الكوكب الذي على طَرَف ذَنَب\textsuperscript{3} الدُّبّ الاصغر & عا & ك & صو & \AbjadZero{} & \multirow{4}{*}{\rotatebox[origin=c]{-90}{\large الشمال}} & ج \\
الذي على اصْل ذَنَب هذا الدبّ الاصغر & قا & ك & عد & ك &  & د \\
الذي من ناحية جَنوب الجَنْب الثاني من المُرَبَّعة & قنج & ك & عب & ى &  & ب \\
الشماليّ من هذا الجَنْب الثاني من المرَبَّعة & قنز & ك & عد & ن &  & ب \\
\hline
\multicolumn{7}{|c|}{\large ومن الدُّبّ الاكبر\textsuperscript{4}} \\ \hline
الكوكب الذي على خَطْم الدُّبّ الاكبر & قعو & ل & لط & كب & \multirow{5}{*}{\rotatebox[origin=c]{-90}{\large الشمال}} & د \\
الذي على رُكْبة هذا الدبّ اليُسْرَى & قيا & ن & له & \AbjadZero{} &  & ج \\
الشماليّ الذي على طَرَف رِجْله اليُسْرَى المقدَّمة & قو & م & كط & ك &  & ج \\
الكوكب الجَنوبيّ منه & قز & ل & كح & ك &  & ج \\
الكوكب الذي على ظَهْره في المُرَبع & قيح & ن & مط & \AbjadZero{} &  & ب \\
الذي على مَرَاقّ بَطْن هذا الدبّ الاكبر & قلج & ك & مد & ى & \multirow{5}{*}{\rotatebox[origin=c]{-90}{\large الشمال}} & ب \\
الكوكب الذي في اصْل ذَنَبه & قلد & \AbjadZero{} & نا & \AbjadZero{} &  & ج \\
الذي على اصل فَخِذه اليُسْرَى المُؤَخَّرة & قلج & ل & مو & ل &  & ب \\
المقدَّم الذي على طَرَف رِجْله اليُسْرَى المُؤَخَّرة & قكج & ل & مط & ك &  & ج \\
الكوكب الذي يتلو هذا & قكه & ك & كح & نه &  & ج \\
\hline
\end{tabular}\end{Arabic}
\par\vspace{3mm}\noindent\rule{40mm}{.4pt}\par\vspace{1mm}{\small 1) Maghrebinice = \textarabic{اقصا} --- 2) Titulus deest in cod. --- 3) Cod. \textarabic{ركبة} --- 4) Titulus deest in cod.}
\clearpage
{\large \textarabic{٢٤٦}}\par\vspace{2mm}
\begin{Arabic}\begin{tabular}{|p{84mm}|>{\centering\arraybackslash}p{9mm}|>{\centering\arraybackslash}p{9mm}|>{\centering\arraybackslash}p{9mm}|>{\centering\arraybackslash}p{9mm}|>{\centering\arraybackslash}p{14mm}|>{\centering\arraybackslash}p{17mm}|}\hline
\multirow{2}{*}{\parbox{80mm}{\centering من اسماء الكواكب الثابتة التي في الصُّوَر الشماليّة من مِنْطَقَة البروج}} & \multicolumn{2}{c|}{الطول} & \multicolumn{2}{c|}{العرض} & \multirow{2}{*}{\parbox{13mm}{\centering\small علامات الجهة}} & \multirow{2}{*}{\parbox{16mm}{\centering\small مراتب العظمة}} \\ \cline{2-5}
 & {\small درج} & {\small دقائق} & {\small درج} & {\small دقائق} & & \\ \hline
الجنوبيّ\textsuperscript{1} من التي على طَرَف رِجْله اليُمْنَى المُؤَخَّرة & قما & ك & مه & مز & \multirow{5}{*}{\rotatebox[origin=c]{-90}{\large الشمال}} & ج \\
الكوكب الجنوبيّ عن هٰؤُلاء & قما & ل & كح & \AbjadZero{} &  & ج \\
المقدَّم من الثلثة على ذَنَبه & قعج & ك & نج & ل &  & ب \\
المُتَوَسِّط من هذه الثلثة & قمط & ى & نه & م &  & ب \\
الثالث الذي على طَرَف ذَنَبه & قصا & \AbjadZero{} & ند & لا &  & ب \\
\hline
\multicolumn{7}{|c|}{\large وممّا ليس\textsuperscript{2} له في صورة الدُّبّ} \\ \hline
الذي تحت الذَّنَب من الجَنوب & قنط & \AbjadZero{} & نط & مه & \multirow{6}{*}{\rotatebox[origin=c]{-90}{\large الشمال}} & ج \\
الذي بين رِجْل الدُّبّ المقدَّمة وبين رأس الأَسَد & قىو & ى & نز & يه &  & د \\
المُظْلِم الذي يتلو الثلثة المُظْلِمة الباقية & قيز & ك & ك & \AbjadZero{} &  & مظلم \\
المُظْلِم المتقدِّم لهذا الكوكب & قيج & ك & كد & ن &  & مظلم \\
الذي بين يَدَيْ هذا الكوكب من المُظْلِمة & قيب & ك & ك & ك &  & مظلم \\
الذي بين يَدَيْه بينه وبين الجَوْزاء & قا & ى & كب & يه &  & مظلم \\
\hline
\multicolumn{7}{|c|}{\large ومن كواكب التِّنِّين\textsuperscript{3}} \\ \hline
الكوكب الذي على طَرَف لِسان التِّنِّين & رىز & ل & عو & ل & \multirow{4}{*}{\rotatebox[origin=c]{-90}{\large الشمال}} & د \\
الكوكب الذي فَوْقَ رأسه & رنا & ى & عه & ى &  & ج \\
الكوكب الذي فوق عَيْن التنّين & رلد & ك & عد & م &  & ج \\
الذي في المغرِب من المُثَلَّث & قصط & ل & فه & ن &  & ج \\
الشماليّ من الاثنَيْن ممّا يَلِي المغرِب & قعا & ى & عح & \AbjadZero{} & \multirow{4}{*}{\rotatebox[origin=c]{-90}{\large الشمال}} & ج \\
الذي يتلو المقدَّم البعيد من الاثنَيْن & قمب & ك & صه & ل &  & ج \\
الذي يتلو هذا الكوكب & قد & ك & صا & مه &  & ج \\
الكوكب الذي على طَرَف ذَنَب التِّنِّين & قد & م & نو & يه &  & ج \\
\hline
\end{tabular}\end{Arabic}
\par\vspace{3mm}\noindent\rule{40mm}{.4pt}\par\vspace{1mm}{\small 1) Ita cod. pro \textarabic{الشماليّ} --- 2) Cod. \textarabic{ليست} --- 3) Inc. fol. 227,r.}
\clearpage
{\large \textarabic{٢٤٧}}\par\vspace{2mm}
\begin{Arabic}\begin{tabular}{|p{84mm}|>{\centering\arraybackslash}p{9mm}|>{\centering\arraybackslash}p{9mm}|>{\centering\arraybackslash}p{9mm}|>{\centering\arraybackslash}p{9mm}|>{\centering\arraybackslash}p{14mm}|>{\centering\arraybackslash}p{17mm}|}\hline
\multirow{2}{*}{\parbox{80mm}{\centering من اسماء الكواكب الثابتة التي في الصُّوَر الشماليّة من مِنْطَقَة البروج}} & \multicolumn{2}{c|}{الطول} & \multicolumn{2}{c|}{العرض} & \multirow{2}{*}{\parbox{13mm}{\centering\small علامات الجهة}} & \multirow{2}{*}{\parbox{16mm}{\centering\small مراتب العظمة}} \\ \cline{2-5}
 & {\small درج} & {\small دقائق} & {\small درج} & {\small دقائق} & & \\ \hline
\multicolumn{7}{|c|}{\large ومن كواكب المُلْتَهِب وهو قيفاوس\textsuperscript{1}} \\ \hline
الكوكب المُضاف الذي على كَتِفه اليُمْنَى & سز & ن & صط & \AbjadZero{} & \multirow{4}{*}{\rotatebox[origin=c]{-90}{\large الشمال}} & ج \\
المُضاف الذي على مِرْفَقه الأَيْمَن & سن & ل & عب & يه &  & د كبير \\
الكوكب الذي على ساعِده الأَيْسَر & صح & يج & صب & ل &  & د \\
المتوسِّط من الثلثة التي على قَلَنْسُوَته & سر & يج & ص & يه &  & د \\
\hline
\multicolumn{7}{|c|}{\large ومن كواكب الغول حارِس الشَّمال وهو البَقَّار\textsuperscript{2}} \\ \hline
الذي على كَتِفه اليُسْرَى & قف & ن & مط & \AbjadZero{} & \multirow{6}{*}{\rotatebox[origin=c]{-90}{\large الشمال}} & ج \\
الكوكب الذي على رأسه & قفز & ن & نج & ن &  & د كبير \\
الكوكب الذي على كَتِفه اليُمْنَى & قضو & ن & مح & م &  & د كبير \\
الكوكب الذي تحت كَتِفه الشَّماليّة & قضح & ن & مو & ل &  & د كبير \\
الذي على فَخِذه اليُمْنَى في المِنْطَقة والرِّباط & قضا & ى & م & يه &  & ج \\
المقدَّم من الاثنَيْن اللذيْن في مِنْطَقته & قفو & ى & مب & ل &  & د كبير \\
الكوكب الذي على عُرْقُوبه الأَيْمَن & قصو & ل & كح & \AbjadZero{} & \multirow{5}{*}{\rotatebox[origin=c]{-90}{\large الشمال}} & ج \\
الشماليّ من الثلثة التي في ساقه اليُسْرَى & قعب & ك & كح & \AbjadZero{} &  & ج \\
الكوكب المُتَوَسِّط من هذه الثلثة & قعا & م & كو & ل &  & د \\
الجَنوبيّ من هذه الثلثة & قعب & ى & كه & \AbjadZero{} &  & ج \\
﴿السِّماك الرامِح﴾ وهو بين فَخِذَيِ الغُول وليس\newline في\textsuperscript{3} صورته & \raisebox{-1\normalbaselineskip}[0pt][0pt]{قفح} & \raisebox{-1\normalbaselineskip}[0pt][0pt]{ى} & \raisebox{-1\normalbaselineskip}[0pt][0pt]{لا} & \raisebox{-1\normalbaselineskip}[0pt][0pt]{ل} &  & \raisebox{-1\normalbaselineskip}[0pt][0pt]{ا} \\
\hline
\end{tabular}\end{Arabic}
\par\vspace{3mm}\noindent\rule{40mm}{.4pt}\par\vspace{1mm}{\small 1) Cod. \textarabic{فيقاوس} --- 2) Cod. \textarabic{النغار} --- 3) Cod. \textarabic{وسرى}}
\clearpage
{\large \textarabic{٢٤٨}}\par\vspace{2mm}
\begin{Arabic}\begin{tabular}{|p{84mm}|>{\centering\arraybackslash}p{9mm}|>{\centering\arraybackslash}p{9mm}|>{\centering\arraybackslash}p{9mm}|>{\centering\arraybackslash}p{9mm}|>{\centering\arraybackslash}p{14mm}|>{\centering\arraybackslash}p{17mm}|}\hline
\multirow{2}{*}{\parbox{80mm}{\centering من اسماء الكواكب الثابتة التي في الصُّوَر الشماليّة من مِنْطَقَة البروج}} & \multicolumn{2}{c|}{الطول} & \multicolumn{2}{c|}{العرض} & \multirow{2}{*}{\parbox{13mm}{\centering\small علامات الجهة}} & \multirow{2}{*}{\parbox{16mm}{\centering\small مراتب العظمة}} \\ \cline{2-5}
 & {\small درج} & {\small دقائق} & {\small درج} & {\small دقائق} & & \\ \hline
\multicolumn{7}{|c|}{\large ومن كواكب الفَكَّة\textsuperscript{1}} \\ \hline
المُنِير من كواكب الفَكَّة & ره & ن & مد & ل & ش & ب كبير \\
المقدَّم من كواكب الفَكَّة & رب & ن & مو & ل & ش & د كبير \\
\hline
\multicolumn{7}{|c|}{\large ومن كواكب الجَاثي} \\ \hline
الذي على رأس الجاثي & رلح & ن & لز & ل & ش & ج \\
الذي على كَتِفه اليُمْنَى عند الإِبْط & ركد & ن & مج & \AbjadZero{} & ش & ج \\
الذي على ذِراعه اليُمْنَى & ركب & ن & م & ى &  & ج \\
الذي على كَتِفه اليُسْرَى & رلز & ن & مح & \AbjadZero{} &  & ج \\
الذي على ذِراعه اليُسْرَى & رمج & ى & مط & ل &  & د كبير \\
الذي على مِرْفَقه الأَيْسَر & رمح & ن & نب & \AbjadZero{} &  & د كبير \\
الجَنوبيّ من الثلثة التي على ساعِده الأَيْسَر & رنب & م & نج & \AbjadZero{} &  & ج \\
الكوكب الذي في الخَطّ الأَيْمَن & ركز & ن & ن & م &  & ج كبير \\
الذي على أَصْل فَخِذه اليُسْرَى & رله & كب & نح & ل &  & ج \\
الذي يتلو هذا في فَخِذه اليُسْرَى من الثلثة & رلز & ل & صا & يه &  & د كبير \\
الذي على اصْل فَخِذه اليُمْنَى & ركج & ى & صد & \AbjadZero{} & ش & د كبير \\
الذي على رُكْبَته اليُمْنَى & رو & ن & صه & ل & ش & د كبير \\
\hline
\multicolumn{7}{|c|}{\large النَّسْر الواقِع} \\ \hline
المضِيء الذي على قَلَنْسُوَة ماسك اللورة\textsuperscript{2} وهو ﴿النَّسْر﴾ & رصح & ل & صب & \AbjadZero{} & ش & ا \\
الشماليّ من الاثنَيْن القريبَيْن منه & رعا & ل & صب & م & ش & د كبير \\
\hline
\end{tabular}\end{Arabic}
\par\vspace{3mm}\noindent\rule{40mm}{.4pt}\par\vspace{1mm}{\small 1) Inc. f. 227,v. --- 2) Cod. \textarabic{اللوزه}}
\clearpage
{\large \textarabic{٢٤٩}}\par\vspace{2mm}
\begin{Arabic}\begin{tabular}{|p{84mm}|>{\centering\arraybackslash}p{9mm}|>{\centering\arraybackslash}p{9mm}|>{\centering\arraybackslash}p{9mm}|>{\centering\arraybackslash}p{9mm}|>{\centering\arraybackslash}p{14mm}|>{\centering\arraybackslash}p{17mm}|}\hline
\multirow{2}{*}{\parbox{80mm}{\centering من اسماء الكواكب الثابتة التي في الصُّوَر الشماليّة من مِنْطَقَة البروج}} & \multicolumn{2}{c|}{الطول} & \multicolumn{2}{c|}{العرض} & \multirow{2}{*}{\parbox{13mm}{\centering\small علامات الجهة}} & \multirow{2}{*}{\parbox{16mm}{\centering\small مراتب العظمة}} \\ \cline{2-5}
 & {\small درج} & {\small دقائق} & {\small درج} & {\small دقائق} & & \\ \hline
الكوكب الجَنُوبيّ منهما & رعا & ل & صا & \AbjadZero{} & ش & د كبير \\
الجنوبيّ من الاثنين اللذيْن في مُقَدَّم كَفَّة الميزان\newline المقدَّمة & \raisebox{-1\normalbaselineskip}[0pt][0pt]{رعب} & \raisebox{-1\normalbaselineskip}[0pt][0pt]{\AbjadZero{}} & \raisebox{-1\normalbaselineskip}[0pt][0pt]{نه} & \raisebox{-1\normalbaselineskip}[0pt][0pt]{\AbjadZero{}} & \raisebox{-1\normalbaselineskip}[0pt][0pt]{ش} & \raisebox{-1\normalbaselineskip}[0pt][0pt]{د صغير} \\
الجنوبيّ من الاثنين اللذيْن في مقدَّم كَفَّة الميزان\newline المُؤَخَّرة & \raisebox{-1\normalbaselineskip}[0pt][0pt]{رعب} & \raisebox{-1\normalbaselineskip}[0pt][0pt]{ى} & \raisebox{-1\normalbaselineskip}[0pt][0pt]{ند} & \raisebox{-1\normalbaselineskip}[0pt][0pt]{مه} &  & \raisebox{-1\normalbaselineskip}[0pt][0pt]{د صغير} \\
الشماليّ الاوّل من الاثنين اللذيْن في كَفَّة الميزان\newline المُؤَخَّرة & \raisebox{-1\normalbaselineskip}[0pt][0pt]{رعب} & \raisebox{-1\normalbaselineskip}[0pt][0pt]{ى} & \raisebox{-1\normalbaselineskip}[0pt][0pt]{نو} & \raisebox{-1\normalbaselineskip}[0pt][0pt]{م} & \raisebox{-1\normalbaselineskip}[0pt][0pt]{ش} & \raisebox{-1\normalbaselineskip}[0pt][0pt]{ج} \\
الشماليّ الثاني من الاثنين اللذيْن في كفّة الميزان\newline المُؤَخَّرة & \raisebox{-1\normalbaselineskip}[0pt][0pt]{رعو} & \raisebox{-1\normalbaselineskip}[0pt][0pt]{ك} & \raisebox{-1\normalbaselineskip}[0pt][0pt]{نه} & \raisebox{-1\normalbaselineskip}[0pt][0pt]{ك} & \raisebox{-1\normalbaselineskip}[0pt][0pt]{ش} & \raisebox{-1\normalbaselineskip}[0pt][0pt]{ج} \\
\hline
\multicolumn{7}{|c|}{\large ومن كواكب الدَّجاجَة\textsuperscript{1}} \\ \hline
الذي على مِنْقار الدَّجاجة & رفز & ك & مط & ك & ش & ج \\
الذي في وَسَط عُنُق الدجاجة & رضز & ل & ند & ل & ش & ج \\
الذي في صَدْر الدجاجة & سط & م & نز & ك &  & ج \\
الكوكب المُضِيء الذي على ذَنَب الدجاجة & سك & ك & ص & \AbjadZero{} &  & ب \\
الكوكب الذي على مِرْفَق جَناحها الأَيْمَن & س & ل & صد & م &  & ج \\
المُتَوَسِّط من الثلثة التي في الجَناح الأَيْسَر\textsuperscript{2} & سب & ك & عد & ى &  & د كبير \\
الكوكب الشماليّ من هذا وهو على طَرَف الجَناح & رضز & ن & عد & \AbjadZero{} &  & د كبير \\
الذي على طَرَف الجناح الأَيْسَر & سيب & \AbjadZero{} & مط & ل &  & ج \\
الذي على رِجْلها اليُسْرَى & سكا & ى & نه & ى & ش & د كبير \\
الذي على رُكْبَته اليُسْرَى & سكو & ك & نز & \AbjadZero{} & ش & د كبير \\
\hline
\end{tabular}\end{Arabic}
\par\vspace{3mm}\noindent\rule{40mm}{.4pt}\par\vspace{1mm}{\small 1) Inc. f. 228,r. --- 2) Codicis error pro \textarabic{الايمن}}
\par\vfill\hfill 32
\clearpage
{\large \textarabic{٢٥٠}}\par\vspace{2mm}
\begin{Arabic}\begin{tabular}{|p{84mm}|>{\centering\arraybackslash}p{9mm}|>{\centering\arraybackslash}p{9mm}|>{\centering\arraybackslash}p{9mm}|>{\centering\arraybackslash}p{9mm}|>{\centering\arraybackslash}p{14mm}|>{\centering\arraybackslash}p{17mm}|}\hline
\multirow{2}{*}{\parbox{80mm}{\centering من اسماء الكواكب الثابتة التي في الصُّوَر الشماليّة من مِنْطَقَة البروج}} & \multicolumn{2}{c|}{الطول} & \multicolumn{2}{c|}{العرض} & \multirow{2}{*}{\parbox{13mm}{\centering\small علامات الجهة}} & \multirow{2}{*}{\parbox{16mm}{\centering\small مراتب العظمة}} \\ \cline{2-5}
 & {\small درج} & {\small دقائق} & {\small درج} & {\small دقائق} & & \\ \hline
\multicolumn{7}{|c|}{\large ومن كواكب ذات الكُرْسِيّ} \\ \hline
الكوكب الذي على رأس ذات الكُرْسِيَ & يط & \AbjadZero{} & مه & ك & ش & ج \\
الكوكب الذي في صَدْرها & نب & \AbjadZero{} & مو & مه & ش & ج \\
الشماليّ منها وهو على شِقَّتها & كج & ن & مز & ن &  & د كبير \\
الكوكب الذي على رُكْبَتها & لا & ن & مه & ل &  & ج \\
الكوكب الذي في وَسَط الكُرْسِيّ & ىط & \AbjadZero{} & نا & م & ش & د صغير \\
الذي في فوْق رِجْل الكرسيّ & كو & ى & نب & م & ش & سَحابِيّ \\
\hline
\multicolumn{7}{|c|}{\large ومن كواكب فرساوس\textsuperscript{1} وهو الفارس المُمْسِك لرأس الغُول} \\ \hline
الذي على طَرَف يد الفارس اليُمْنَى وهو مُمْسِك\newline رَأْس الغُول & \raisebox{-1\normalbaselineskip}[0pt][0pt]{لز} & \raisebox{-1\normalbaselineskip}[0pt][0pt]{ن} & \raisebox{-1\normalbaselineskip}[0pt][0pt]{م} & \raisebox{-1\normalbaselineskip}[0pt][0pt]{ل} & \raisebox{-1\normalbaselineskip}[0pt][0pt]{ش} & \raisebox{-1\normalbaselineskip}[0pt][0pt]{د صغير} \\
المُضِيء الذي في شِقّه الأَيمن & مو & \AbjadZero{} & ل & \AbjadZero{} & ش & ج \\
الذي على كَتِفه اليُمْنَى & مج & ن & لد & ل &  & ب \\
المُؤَخَّر من الثلثة التي في شِقّه الأَيمن & مح & ن & كز & ك &  & د كبير \\
الذي على فَخِذه اليُسْرَى & مج & \AbjadZero{} & كا & ن &  & د \\
المُنِير من التي في رأس الغول & م & ن & كج & \AbjadZero{} &  & ب \\
الكوكب\textsuperscript{2} الذي على عُرْقُوبه الأَيْسَر & مه & \AbjadZero{} & يب & \AbjadZero{} & ش & ج صغير \\
الذي يتلو هذا وهو على قَدَمه اليُسْرَى & مز & ل & يا & \AbjadZero{} & ش & ج \\
\hline
\end{tabular}\end{Arabic}
\par\vspace{3mm}\noindent\rule{40mm}{.4pt}\par\vspace{1mm}{\small 1) Cod. \textarabic{فيقاوس} --- 2) Inc. f. 228,v.}
\clearpage
{\large \textarabic{٢٥١}}\par\vspace{2mm}
\begin{Arabic}\begin{tabular}{|p{84mm}|>{\centering\arraybackslash}p{9mm}|>{\centering\arraybackslash}p{9mm}|>{\centering\arraybackslash}p{9mm}|>{\centering\arraybackslash}p{9mm}|>{\centering\arraybackslash}p{14mm}|>{\centering\arraybackslash}p{17mm}|}\hline
\multirow{2}{*}{\parbox{80mm}{\centering من اسماء الكواكب الثابتة التي في الصُّوَر الشماليّة من مِنْطَقَة البروج}} & \multicolumn{2}{c|}{الطول} & \multicolumn{2}{c|}{العرض} & \multirow{2}{*}{\parbox{13mm}{\centering\small علامات الجهة}} & \multirow{2}{*}{\parbox{16mm}{\centering\small مراتب العظمة}} \\ \cline{2-5}
 & {\small درج} & {\small دقائق} & {\small درج} & {\small دقائق} & & \\ \hline
\multicolumn{7}{|c|}{\large ومن كواكب ذي الأَعِنَّة وهو العَنَّاز\textsuperscript{1} الجَنُوبيّ} \\ \hline
الكوكب الذي على رأْس ذي الأَعِنَّة & عج & م & ل & \AbjadZero{} & ش & ج \\
الكوكب الذي في كَتِفه اليُسْرَى وهو ﴿العَيُّوق﴾\textsuperscript{2} & صو & ى & كب & ل & ش & ا\textsuperscript{3} \\
الذي على كَتِفه اليُمْنَى & عد & \AbjadZero{} & ك & \AbjadZero{} &  & ب \\
الكوكب الذي على عُرْقوبه الأَيْسَر & صا & \AbjadZero{} & ى & ى & ش & ج صغير \\
الذي على عُرْقوبه الأَيْمَن & صر & ى & ه & \AbjadZero{} & ش & ج كبير \\
\hline
\multicolumn{7}{|c|}{\large ومن كواكب الحَوَّاء\textsuperscript{4} الذي يُمْسِك الحَيَّة} \\ \hline
الكوكب الذي على رأس الحَوَّاء\textsuperscript{5} وهو مُمْسِك الحَيَّة & رمو & \AbjadZero{} & لو & \AbjadZero{} & ش & ج كبير \\
المقدَّم من الاثنين اللذيْن في كَتِفه اليُمْنَى & رمط & ى & كز & يه & ش & د \\
الاوّل من الاثنين اللذيْن على طَرَف يده اليُسْرَى & ركو & ى & يز & \AbjadZero{} &  & ج \\
الكوكب الذي يتلو هذا في طَرَف اليد & ركز & ى & يو & ل &  & ج \\
الكوكب الذي على رُكْبَته اليُمْنَى & رمب & م & ز & ل &  & ج \\
الكوكب الذي على ساقه اليُمْنَى\textsuperscript{6} & رمد & ن & ب & يه &  & ج كبير \\
الثاني من الاربعة التي على رِجْله & رمه & ل & ا & له & ش & د كبير \\
الشماليّ الذي في خطّ الثلثة الشماليّة & رنج & ك & يا & ن & ش & ج كبير \\
\hline
\multicolumn{7}{|c|}{\large ومن كواكب الحَيَّة التي يُمْسِكها الحَوَّاء\textsuperscript{7}} \\ \hline
الكوكب الذي في صُدْغ الحَيَّة التي يُمْسِكها الحَوَّاء\textsuperscript{8} & رنه & ل & لو & \AbjadZero{} & ش & ج \\
المُضاف الذي في مِنْخَرَيِ الحَيَّة & رنب & ن & م & \AbjadZero{} & ش & د \\
\hline
\end{tabular}\end{Arabic}
\par\vspace{3mm}\noindent\rule{40mm}{.4pt}\par\vspace{1mm}{\small 1) Cod. \textarabic{العنان} --- 2) Cod. \textarabic{العير} --- 3) Cod. addit \textarabic{سحابى} --- 4) Cod. \textarabic{الجوا} --- 5) Cod. \textarabic{الجوا} --- 6) Cod. \textarabic{الايمن} --- 7) Cod. \textarabic{الجوى} --- 8) Cod. \textarabic{الجوا}}
\clearpage
{\large \textarabic{٢٥٢}}\par\vspace{2mm}
\begin{Arabic}\begin{tabular}{|p{84mm}|>{\centering\arraybackslash}p{9mm}|>{\centering\arraybackslash}p{9mm}|>{\centering\arraybackslash}p{9mm}|>{\centering\arraybackslash}p{9mm}|>{\centering\arraybackslash}p{14mm}|>{\centering\arraybackslash}p{17mm}|}\hline
\multirow{2}{*}{\parbox{80mm}{\centering من اسماء الكواكب الثابتة التي في الصُّوَر الشماليّة من مِنْطَقَة البروج}} & \multicolumn{2}{c|}{الطول} & \multicolumn{2}{c|}{العرض} & \multirow{2}{*}{\parbox{13mm}{\centering\small علامات الجهة}} & \multirow{2}{*}{\parbox{16mm}{\centering\small مراتب العظمة}} \\ \cline{2-5}
 & {\small درج} & {\small دقائق} & {\small درج} & {\small دقائق} & & \\ \hline
الذي عند مَخْرَج عُنُق الحَيّة & رنج & ى & لد & ى & ش & ج \\
الذي بعد الخَرَزَة\textsuperscript{1} المتقدِّمة التي في عُنُق الحَيّة & رنب & ن & نط & يه & ش & ج \\
المتوسِّط من الثلثة التي بعدها & رنه & ل & كه & ك &  & ج \\
الكوكب الجنوبيّ منها & رنز & ل & كد & \AbjadZero{} &  & ج كبير \\
الكوكب\textsuperscript{2} الجنوبيّ الذي وَرَاءَ فَخِذ الحَوَّاء\textsuperscript{3} & رنح & ى & يز & ل &  & د كبير \\
الذي يَلِي الاوّل من الثلثة التي على ذَنَب الحَيّة & رنط & ن & كا & يه & ش & د كبير \\
الذي على طَرَف ذَنَب الحَيَّة & رصط & ل & كز & \AbjadZero{} & ش & د \\
\hline
\multicolumn{7}{|c|}{\large من كواكب اويسطس\textsuperscript{4} وهو النَّصْل\textsuperscript{5}} \\ \hline
الكوكب الفَرِيد الذي على سَهْمه & رضا & ك & لط & ك & ش & د \\
الكوكب الذي على آخِر السَّهْم & رفد & ل & لح & م & ش & د \\
\hline
\multicolumn{7}{|c|}{\large ومن كواكب النَّسْر الطَّائِر} \\ \hline
المُقَدَّم من الاثنَيْن اللذيْن في كَتِفه اليُسْرَى\textsuperscript{6} & رضد & ك & لا & ل & ش & ج \\
الذي يتلو الذي في وَسَط رأسه وهو على عُنُقه & رفد & ك & كز & ى & ش & ج \\
﴿النَّسْر الطائر﴾ وهو المُضِيء الذي بين كَتِفَيْه & رفه & \AbjadZero{} & كط & ى &  & ج كبير \\
الكوكب الشماليّ القريب من النسر الطائر & رفو & ى & ل & \AbjadZero{} & ش & ج صغير \\
الذي تحت النسر وهو بعيد عنه في لِزْق المَجَرَّة\textsuperscript{7} & رعج & ك & لو & ك & ش & ج \\
\hline
\end{tabular}\end{Arabic}
\par\vspace{3mm}\noindent\rule{40mm}{.4pt}\par\vspace{1mm}{\small 1) Cod. \textarabic{الجرزه} --- 2) Inc. f. 229,r. --- 3) Cod. \textarabic{الجوا} --- 4) Cod. \textarabic{اوسطيس} --- 5) Ex coniectura; cod. \textarabic{النول} --- 6) Cod. \textarabic{الايسر} --- 7) Cod. \textarabic{المجضرة}}
\clearpage
{\large \textarabic{٢٥٣}}\par\vspace{2mm}
\begin{Arabic}\begin{tabular}{|p{84mm}|>{\centering\arraybackslash}p{9mm}|>{\centering\arraybackslash}p{9mm}|>{\centering\arraybackslash}p{9mm}|>{\centering\arraybackslash}p{9mm}|>{\centering\arraybackslash}p{14mm}|>{\centering\arraybackslash}p{17mm}|}\hline
\multirow{2}{*}{\parbox{80mm}{\centering من اسماء الكواكب الثابتة التي في الصُّوَر الشماليّة من مِنْطَقَة البروج}} & \multicolumn{2}{c|}{الطول} & \multicolumn{2}{c|}{العرض} & \multirow{2}{*}{\parbox{13mm}{\centering\small علامات الجهة}} & \multirow{2}{*}{\parbox{16mm}{\centering\small مراتب العظمة}} \\ \cline{2-5}
 & {\small درج} & {\small دقائق} & {\small درج} & {\small دقائق} & & \\ \hline
\multicolumn{7}{|c|}{\large ومن كواكب الدُّلْفِين وهو الصَّلِيب} \\ \hline
المقدَّم من الثلثة التي على ذَنَب الدُّلْفين وهو الصَّلِيب & رضح & ن & كط & ى & ش & ج صغير \\
الجَنوبيّ المقدَّم من الضِّلْع الاوّل & رضط & م & لب & \AbjadZero{} & ش & ج صغير \\
المُضاف الجنوبيّ الذي في خَطّ سَنَد المِنْقار & سب & ى & لج & \AbjadZero{} &  & ج صغير \\
الشماليّ من الضِّلْع المُؤَخَّر & سد & م & لج & ى & ش & ج صغير \\
الشماليّ من الضلع الاوَّل & سا & ك & لج & ن & ش & ج صغير \\
\hline
\multicolumn{7}{|c|}{\large ومن كواكب برطومس وهو الفَرَس} \\ \hline
الاوّل المقدَّم من اللذَيْن في رَأْس بُرُطُومِس وهو\newline الفَرَس & \raisebox{-1\normalbaselineskip}[0pt][0pt]{سز} & \raisebox{-1\normalbaselineskip}[0pt][0pt]{ل} & \raisebox{-1\normalbaselineskip}[0pt][0pt]{ك} & \raisebox{-1\normalbaselineskip}[0pt][0pt]{ل} & \raisebox{-1\normalbaselineskip}[0pt][0pt]{ش} & \raisebox{-1\normalbaselineskip}[0pt][0pt]{مُظْلِم} \\
الثاني وهو المُؤَخَّر منهما & سط & ى & كه & م & ش & مظلم \\
المقدَّم من الاثنَيْن اللذَيْن في فَمِه & سز & ل & كه & ل &  & مظلم \\
الكوكب المُؤَخَّر منهما & سح & \AbjadZero{} & كه & ل &  & مظلم \\
المُتَوَسِّط من التي في رأْس اندروميذس\textsuperscript{1} وهي المَرْأَة\newline التي ليس لها بَعْل & \raisebox{-1\normalbaselineskip}[0pt][0pt]{سٮح} & \raisebox{-1\normalbaselineskip}[0pt][0pt]{ل} & \raisebox{-1\normalbaselineskip}[0pt][0pt]{كو} & \raisebox{-1\normalbaselineskip}[0pt][0pt]{\AbjadZero{}} &  & \raisebox{-1\normalbaselineskip}[0pt][0pt]{ب صغير} \\
الكوكب\textsuperscript{2} الذي على ظَهْر الفَرَس وهو في رأس كَتِفه & سنج & ك & يب & ل &  & ب صغير \\
﴿مَنْكِب الفَرَس﴾ وهو على كَتِفه اليُمْنَى في مَخْرَج\newline قدَم\textsuperscript{3} الفَرَس & \raisebox{-1\normalbaselineskip}[0pt][0pt]{سلج} & \raisebox{-1\normalbaselineskip}[0pt][0pt]{ك} & \raisebox{-1\normalbaselineskip}[0pt][0pt]{لا} & \raisebox{-1\normalbaselineskip}[0pt][0pt]{\AbjadZero{}} &  & \raisebox{-1\normalbaselineskip}[0pt][0pt]{ب صغير} \\
الذي بين كَتِفَيْه في كَتِف جَناح الفَرَس & سلز & ن & يط & م &  & ب صغير \\
الشماليّ من الاثنَيْن اللذيْن في رُكْبَته اليُمْنَى & سم & ى & له & \AbjadZero{} & ش & ج \\
المقدَّم من الاثنَيْن اللذيْن في عُنُقه & سل & \AbjadZero{} & يح & \AbjadZero{} & ش & ج \\
\hline
\end{tabular}\end{Arabic}
\par\vspace{3mm}\noindent\rule{40mm}{.4pt}\par\vspace{1mm}{\small 1) Cod. \textarabic{اندرومذرس} --- 2) Inc. f. 229,v. --- 3) Deest in cod.}
\clearpage
{\large \textarabic{٢٥٤}}\par\vspace{2mm}
\begin{Arabic}\begin{tabular}{|p{84mm}|>{\centering\arraybackslash}p{9mm}|>{\centering\arraybackslash}p{9mm}|>{\centering\arraybackslash}p{9mm}|>{\centering\arraybackslash}p{9mm}|>{\centering\arraybackslash}p{14mm}|>{\centering\arraybackslash}p{17mm}|}\hline
\multirow{2}{*}{\parbox{80mm}{\centering من اسماء الكواكب الثابتة التي في الصُّوَر الشماليّة من مِنْطَقَة البروج}} & \multicolumn{2}{c|}{الطول} & \multicolumn{2}{c|}{العرض} & \multirow{2}{*}{\parbox{13mm}{\centering\small علامات الجهة}} & \multirow{2}{*}{\parbox{16mm}{\centering\small مراتب العظمة}} \\ \cline{2-5}
 & {\small درج} & {\small دقائق} & {\small درج} & {\small دقائق} & & \\ \hline
الشماليّ من الاثنَيْن اللذيْن في رأسه & سك & ك & نو & ن & ش & ج \\
الكوكب الذي على حُقُومه & سمو & لج & كب & ل & ش & ج كبير \\
الذي على عُرْقوبه الأَيْمَن & سلد & ن & كا & ى &  & د كبير \\
الكوكب الذي على رُكْبَته اليُسْرَى & سكح & ل & لد & يه & ش & د كبير \\
الكوكب الذي على عُرْقوبه الأَيْسَر & سكج & ل & لو & ن & ش & د كبير \\
\hline
\multicolumn{7}{|c|}{\large ومن كواكب اندروميذس\textsuperscript{1} وهي المَرْأَة التي لَمْ تَرَ بَعْلًا} \\ \hline
الذي بين كَتِفَيْ اندروميذس\textsuperscript{2} & و & ل & كد & ل & ش & ج \\
الجنوبيّ من الثلثة التي فوْق شِقَّتها\textsuperscript{3} & يه & \AbjadZero{} & كو & ك & ش & ج \\
المقدَّم الخارج من الثلثة التي في رأسها\textsuperscript{4} & سله & ن & مد & \AbjadZero{} & ش & ج \\
الكوكب الذي فوْق رِجْلها اليُسْرَى & كح & \AbjadZero{} & كح & \AbjadZero{} & ش & ج \\
\hline
\multicolumn{7}{|c|}{\large ومن كواكب طريغانس\textsuperscript{5} وهو المُثَلَّث} \\ \hline
الكوكب الذي في رَأْس المُثَلَّث & فب & ى & يو & ل & ش & ج \\
المقدَّم من الثلثة التي في أَسْفَله & كج & ى & ك & م &  & ج \\
الكوكب الذي في آخِر هذه الثلثة & كح & \AbjadZero{} & يط & \AbjadZero{} & ش & ج \\
\hline
\multicolumn{7}{|c|}{\parbox{168mm}{\centering\large ابتداء\textsuperscript{6} أسْماء الكواكب الثابتة التي في الصُّوَر الظاهرة في مِنْطَقَة البروج لسنة \AbjadOver{اقضا}\textsuperscript{7} لذي\\القَرْنَيْن وهي السنة \AbjadOver{رعا} للهِجْرَة}} \\ \hline
\multicolumn{7}{|c|}{\large اسماء كواكب صورة الحَمَل} \\ \hline
الكوكب المقدَّم من الاثنَيْن اللذيْن في قَرْن الحَمَل & يد & م & ز & ك & ش & ج صغير \\
\hline
\end{tabular}\end{Arabic}
\par\vspace{3mm}\noindent\rule{40mm}{.4pt}\par\vspace{1mm}{\small 1) Cod. \textarabic{اندرمندش} --- 2) Cod. \textarabic{اندرومندس} --- 3) Cod. \textarabic{شفتها} --- 4) Error pro \textarabic{كَفَّها اليُمْنَى} --- 5) Cod. \textarabic{طربغاش} --- 6) Inc. fol. 230,r. --- 7) Maghrebinice = \textarabic{اقصا} Orientalium.}
\clearpage
{\large \textarabic{٢٥٥}}\par\vspace{2mm}
\begin{Arabic}\begin{tabular}{|p{84mm}|>{\centering\arraybackslash}p{9mm}|>{\centering\arraybackslash}p{9mm}|>{\centering\arraybackslash}p{9mm}|>{\centering\arraybackslash}p{9mm}|>{\centering\arraybackslash}p{14mm}|>{\centering\arraybackslash}p{17mm}|}\hline
\multirow{2}{*}{\parbox{80mm}{\centering من اسماء الكواكب الثابتة التي في صُوَر مِنْطَقَة البروج}} & \multicolumn{2}{c|}{الطول} & \multicolumn{2}{c|}{العرض} & \multirow{2}{*}{\parbox{13mm}{\centering\small علامات الجهة}} & \multirow{2}{*}{\parbox{16mm}{\centering\small مراتب العظمة}} \\ \cline{2-5}
 & {\small درج} & {\small دقائق} & {\small درج} & {\small دقائق} & & \\ \hline
الكوكب المُؤَخَّر منهما & يح & م & ح & ك & ش & ج \\
الشماليّ من الاثنَيْن اللذَيْنِ في فَم الحَمَل & كب & ى & ز & م & ش & ه \\
الكوكب الجنوبيّ منهما & كب & ك & و & \AbjadZero{} &  & ه \\
الكوكب الذي على عُنُق الحَمَل & يز & م & ه & ل &  & ه \\
الذي على ظَهْر الحمل & كح & م & و & \AbjadZero{} &  & و \\
الذي على مَخْرَج أَلْيَته & لب & ل & د & م &  & ه \\
المقدَّم من الثلثة التي على ألْيته & له & \AbjadZero{} & ا & ن &  & د \\
الكوكب المُتَوَسِّط منها & لو & ل & ب & ل &  & د \\
المُؤَخَّر من هذه الثلثة & لح & ى & ا & ن & ش & د \\
الكوكب الذي خَلْفَ فَخِذ الحَمَل & ل & ل & ا & ى & ش & ه \\
الذي فوْق وَسَط فَخِذه & كط & ى & ا & ى & ج & ه \\
الذي على ظَهْر رِجْله المُؤَخَّرة & كو & ى & ه & يه & ج & د كبير \\
\hline
\multicolumn{7}{|c|}{\large وممّا لَيْسَ في صورته} \\ \hline
وهو الذي فَوْقَ رَأْسه & كا & ل & ى & \AbjadZero{} & ش & ج \\
\hline
\multicolumn{7}{|c|}{\large من أسْماء الكواكب في صورة الثوْر} \\ \hline
الشَّماليّ من الاربعة التي على قَطْع\textsuperscript{1} الثَّوْر & لو & ل & و & \AbjadZero{} & ج & د \\
الثاني الذي يتلوه من هذه الاربعة & لو & ن & ز & يه & ج & د \\
الجَنوبيّ من هذه الاربعة التي على قَطْعه\textsuperscript{2} & له & ل & ط & يه &  & د \\
الكوكب الذي على صَدْر الثَّوْر & مو & ل & ح & \AbjadZero{} &  & ج \\
الذي على كَتِفه اليُمْنَى & م & ل & ط & ن & ج & د \\
الذي في وَجْه الٮوْر على أَنْفه من كواكب الدَّبَران & ن & \AbjadZero{} & ه & مه & ج & ج صغير \\
\hline
\end{tabular}\end{Arabic}
\par\vspace{3mm}\noindent\rule{40mm}{.4pt}\par\vspace{1mm}{\small 1) Cod. \textarabic{قطن} --- 2) Cod. \textarabic{قطنه}}
\clearpage
{\large \textarabic{٢٥٦}}\par\vspace{2mm}
\begin{Arabic}\begin{tabular}{|p{84mm}|>{\centering\arraybackslash}p{9mm}|>{\centering\arraybackslash}p{9mm}|>{\centering\arraybackslash}p{9mm}|>{\centering\arraybackslash}p{9mm}|>{\centering\arraybackslash}p{14mm}|>{\centering\arraybackslash}p{17mm}|}\hline
\multirow{2}{*}{\parbox{80mm}{\centering من اسماء الكواكب الثابتة التي في صُوَر مِنْطَقَة البروج}} & \multicolumn{2}{c|}{الطول} & \multicolumn{2}{c|}{العرض} & \multirow{2}{*}{\parbox{13mm}{\centering\small علامات الجهة}} & \multirow{2}{*}{\parbox{16mm}{\centering\small مراتب العظمة}} \\ \cline{2-5}
 & {\small درج} & {\small دقائق} & {\small درج} & {\small دقائق} & & \\ \hline
الذي\textsuperscript{1} بين هذا الكوكب وبين عَيْنه الشَّماليَّة & نا & ل & ا & يه & ج & ج صغير \\
﴿الدَّبَران﴾ وهو الذي على عَيْنه وتحت قَرْنه الجنوبيّ & نج & ن & ه & ى & ج & ا \\
الذي بين الاوَّل وبين عينه الأُخْرَى الجنوبيّة & نب & \AbjadZero{} & ه & ن &  & ج صغير \\
التالي لهذا الذي على عينه الشماليّة & نج & \AbjadZero{} & ج & \AbjadZero{} &  & د \\
الذي على أَصْل قَرْنه عند أُذْنه الجنوبيّة & نز & م & د & \AbjadZero{} &  & ه \\
الذي على وَسَط قَرْنه الجنوبيّ من الاثنَيْن & صا & ل & ه & \AbjadZero{} &  & ه \\
الذي على اصْل قَرْنه الشماليّ & نو & ن & د & \AbjadZero{} & ج & د \\
الذي على طَرَف قَرْنه الجنوبيّ & صح & ن & ب & ن & ج & ج \\
الذي على طَرَف قَرْنه الشماليّ وهو على رِجْل ذي\newline الأَعِنَّة & \raisebox{-1\normalbaselineskip}[0pt][0pt]{صو} & \raisebox{-1\normalbaselineskip}[0pt][0pt]{ن} & \raisebox{-1\normalbaselineskip}[0pt][0pt]{ه} & \raisebox{-1\normalbaselineskip}[0pt][0pt]{\AbjadZero{}} & \raisebox{-1\normalbaselineskip}[0pt][0pt]{ش} & \raisebox{-1\normalbaselineskip}[0pt][0pt]{ج\textsuperscript{2}} \\
الشماليّ الذي في السَّطْر المقدَّم من الثُّرَيَّا & مج & ك & ه & ل & ش & ج\textsuperscript{2} \\
الشماليّ الذي في آخِر الشِّقّ المقدَّم من الثُّرَيَّا & مج & ل & ج & م &  & د\textsuperscript{2} \\
المُؤَخَّر الصغير الذي مُؤَخَّر الثريّا & مد & ن & ج & ك & ش & ج\textsuperscript{2} \\
الصغير الخارج من شمال الثرَيَّا & مد & ن & ه & \AbjadZero{} & ش & د\textsuperscript{2} \\
\hline
\multicolumn{7}{|c|}{\large أسْماء الكواكب التي في صورتَيِ التَّوْءَمَيْن} \\ \hline
الذي على رَأْس التَّوْءَم المقدَّم & ضد & ل & ط & ك & ش & ب \\
الذي على رأْس التَّوْءَم المُؤَخَّر & ضز & ن & و & يه & ش & ب \\
الذي بين يَدَيْ رِجْل التوْءَم المقدَّم & عز & م & ا & ل & ج & د كبير \\
الذي على طَرَف رِجْل التوْءَم المقدَّم & قا & ك & ج & ل & ج & د كبير \\
الذي على طَرَف الرِّجْل اليُسْرَى من التوْءَم المُؤَخَّر & قج & ى & ز & ل & ج & ج \\
\hline
\end{tabular}\end{Arabic}
\par\vspace{3mm}\noindent\rule{40mm}{.4pt}\par\vspace{1mm}{\small 1) Inc. f. 230,v. --- 2) Cod. perperam addit \textarabic{سحابى}}
\clearpage
{\large \textarabic{٢٥٧}}\par\vspace{2mm}
\begin{Arabic}\begin{tabular}{|p{84mm}|>{\centering\arraybackslash}p{9mm}|>{\centering\arraybackslash}p{9mm}|>{\centering\arraybackslash}p{9mm}|>{\centering\arraybackslash}p{9mm}|>{\centering\arraybackslash}p{14mm}|>{\centering\arraybackslash}p{17mm}|}\hline
\multirow{2}{*}{\parbox{80mm}{\centering من اسماء الكواكب الثابتة التي في صُوَر مِنْطَقَة البروج}} & \multicolumn{2}{c|}{الطول} & \multicolumn{2}{c|}{العرض} & \multirow{2}{*}{\parbox{13mm}{\centering\small علامات الجهة}} & \multirow{2}{*}{\parbox{16mm}{\centering\small مراتب العظمة}} \\ \cline{2-5}
 & {\small درج} & {\small دقائق} & {\small درج} & {\small دقائق} & & \\ \hline
الذي على طرف الرجل اليُمْنَى من التوءم المُؤَخَّر & قه & ن & ى & ل & ج & د \\
الذي على الرُّكْبة اليُسْرَى من التوءم المُؤَخَّر & قط & م & ب & ل & ج & ج \\
الذي على وَسَط اللَّيَّة\textsuperscript{1} اليُمْنَى من هذا التوْءَم & ضب & ن & و & \AbjadZero{} &  & ج \\
الذي على الخاصِرة اليُسْرَى من هذا التوْءَم & ضب & ن & \AbjadZero{} & ل & ج & ج \\
الذي على الرُّكْبة اليُسْرَى من التوْءَم المقدَّم & قد & ى & ا & ى & ج & د كبير \\
\hline
\multicolumn{7}{|c|}{\large اسماء\textsuperscript{2} الكواكب التي في صورة السَّرَطان} \\ \hline
وَسَط المِعْلَف وهو السَّحابيّ الذي في صَدْر السَّرَطان & قيا & ل & \AbjadZero{} & م & ش & سَحابيّ \\
الشماليّ من الاثنَيْن اللذيْن عند المُرَبَّع السَّحابيّ في\newline الإظْلام & \raisebox{-1\normalbaselineskip}[0pt][0pt]{قيا} & \raisebox{-1\normalbaselineskip}[0pt][0pt]{ن} & \raisebox{-1\normalbaselineskip}[0pt][0pt]{ا} & \raisebox{-1\normalbaselineskip}[0pt][0pt]{يه} & \raisebox{-1\normalbaselineskip}[0pt][0pt]{ش} & \raisebox{-1\normalbaselineskip}[0pt][0pt]{د صغير} \\
الكوكب الجنوبيّ منهما ويُسَمَّيانِ ﴿المُرَبَّع﴾ & قيب & ى & يا & ى & ج & د صغير \\
الذي على زُبانة السَّرَطان الشماليّة & قط & ل & يا & ن & ش & د كبير \\
الذي على زُبانة السَّرَطان الجنوبيّة & قيا & م & ه & ل & ج & د \\
الذي على طَرَف رِجْله المُؤَخَّرة الجنوبيّة & قح & ك & ز & ل & ش & ه \\
الذي على طَرَف رِجْله المُؤَخَّرة الشماليّة & قج & ن & ا & \AbjadZero{} & ش & د \\
الجنوبيّ من الاثنَيْن اللذيْن ذكرناهما & قح & ك & ز & ل & ج & د كبير \\
الشماليّ من الاثنين اللذيْن عند السَّحابيّ وهما\newline ﴿الحِمَارَانِ﴾\textsuperscript{1} & \raisebox{-1\normalbaselineskip}[0pt][0pt]{قيا} & \raisebox{-1\normalbaselineskip}[0pt][0pt]{ل} & \raisebox{-1\normalbaselineskip}[0pt][0pt]{ب} & \raisebox{-1\normalbaselineskip}[0pt][0pt]{م} & \raisebox{-1\normalbaselineskip}[0pt][0pt]{ش} & \raisebox{-1\normalbaselineskip}[0pt][0pt]{د كبير} \\
الكوكب الجنوبيّ منهما & قيب & ل & ه & ى & ج & د \\
\hline
\end{tabular}\end{Arabic}
\par\vspace{3mm}\noindent\rule{40mm}{.4pt}\par\vspace{1mm}{\small 1) Cod. \textarabic{اللحه} --- 2) Inc. f. 231,r. --- 3) Cod. \textarabic{وهو الجادان}}
\par\vfill\hfill 33
\clearpage
{\large \textarabic{٢٥٨}}\par\vspace{2mm}
\begin{Arabic}\begin{tabular}{|p{84mm}|>{\centering\arraybackslash}p{9mm}|>{\centering\arraybackslash}p{9mm}|>{\centering\arraybackslash}p{9mm}|>{\centering\arraybackslash}p{9mm}|>{\centering\arraybackslash}p{14mm}|>{\centering\arraybackslash}p{17mm}|}\hline
\multirow{2}{*}{\parbox{80mm}{\centering من اسماء الكواكب الثابتة التي في صُوَر مِنْطَقَة البروج}} & \multicolumn{2}{c|}{الطول} & \multicolumn{2}{c|}{العرض} & \multirow{2}{*}{\parbox{13mm}{\centering\small علامات الجهة}} & \multirow{2}{*}{\parbox{16mm}{\centering\small مراتب العظمة}} \\ \cline{2-5}
 & {\small درج} & {\small دقائق} & {\small درج} & {\small دقائق} & & \\ \hline
\multicolumn{7}{|c|}{\large اسماء كواكب صورة الأَسَد} \\ \hline
الكوكب الذي على طَرَف مِنْخَر الأَسَد & قيط & ل & ى & \AbjadZero{} & ش & د \\
الشماليّ من الاثنَيْن اللذَيْن في رأْس الاسد & قكه & \AbjadZero{} & ىب & \AbjadZero{} & ش & ج \\
الأَوْسَط من الثلثة التي في عُنُقه & قلج & ك & ح & ل &  & ب \\
المقدَّم من الثلثة التي في عُنُقه وهو الشماليّ & قله & ك & يا & \AbjadZero{} &  & د \\
الجنوبيّ من هذه الثلثة التي في عُنُق الاسد & قا & ن & د & ل &  & ج \\
﴿قَلْب الأَسَد﴾ يُسَمَّى ﴿الماكن﴾\textsuperscript{1} & قلد & \AbjadZero{} & \AbjadZero{} & ى &  & ا \\
المقدَّم من الاثنَيْن اللذَيْن على ظَهْره & قمب & ل & يب & يه &  & ه \\
الكوكب الذي يتلو هذا & قمد & ك & يج & م &  & د صغير \\
الذي في فَخِذ الاسد المُؤَخَّرة & قنب & ل & د & ن & ش & ج \\
الثاني الذي في فَخِذه المُؤَخَّرة ايضًا & قنب & ن & ا & يه & ش & د \\
الذي في وَسَط فَخِذه المؤّخرة ايضًا & قنح & م & ج & \AbjadZero{} & ج & ه \\
﴿الصَّرْفَة﴾ وهو الكوكب الذي على طَرَف ذَنَبه & قند & م & يا & ل & ش & ا \\
الشماليّ\textsuperscript{2} من الاثنيْن اللذَيْن في مُؤَخَّر الاسد & قمه & ل & يا & ل & ش & ه \\
الكوكب الجنوبيّ منهما & قمز & ل & ط & م & ش & ج \\
الكوكب الذي على رُكْبة الاسد اليُمْنَى & قكح & ك & \AbjadZero{} & \AbjadZero{} & \AbjadZero{} & ه \\
\hline
\multicolumn{7}{|c|}{\large من كواكب الذُّؤَابة ولَيْسَتْ من صورة الأَسَد} \\ \hline
اوّلها ﴿بُلُوقامُس﴾\textsuperscript{3} وهو الكوكب الذي بين\newline ذَنب الاسد والسِّماك الرامِح & \raisebox{-1\normalbaselineskip}[0pt][0pt]{قنو} & \raisebox{-1\normalbaselineskip}[0pt][0pt]{\AbjadZero{}} & \raisebox{-1\normalbaselineskip}[0pt][0pt]{ل} & \raisebox{-1\normalbaselineskip}[0pt][0pt]{يه} & \raisebox{-1\normalbaselineskip}[0pt][0pt]{ش} & \raisebox{-1\normalbaselineskip}[0pt][0pt]{مشهور} \\
\hline
\end{tabular}\end{Arabic}
\par\vspace{3mm}\noindent\rule{40mm}{.4pt}\par\vspace{1mm}{\small 1) Error videtur pro \textarabic{المَلَكيّ} --- 2) Incipit f. 231,v. --- 3) Cod. \textarabic{لطرقانش}}
\clearpage
{\large \textarabic{٢٥٩}}\par\vspace{2mm}
\begin{Arabic}\begin{tabular}{|p{84mm}|>{\centering\arraybackslash}p{9mm}|>{\centering\arraybackslash}p{9mm}|>{\centering\arraybackslash}p{9mm}|>{\centering\arraybackslash}p{9mm}|>{\centering\arraybackslash}p{14mm}|>{\centering\arraybackslash}p{17mm}|}\hline
\multirow{2}{*}{\parbox{80mm}{\centering من اسماء الكواكب الثابتة التي في صُوَر مِنْطَقَة البروج}} & \multicolumn{2}{c|}{الطول} & \multicolumn{2}{c|}{العرض} & \multirow{2}{*}{\parbox{13mm}{\centering\small علامات الجهة}} & \multirow{2}{*}{\parbox{16mm}{\centering\small مراتب العظمة}} \\ \cline{2-5}
 & {\small درج} & {\small دقائق} & {\small درج} & {\small دقائق} & & \\ \hline
المقدَّم الكبير الذي على الضَّفِيرة\textsuperscript{1} ويُسَمَّى ﴿عُرْف\newline الأَسَد﴾ & \raisebox{-1\normalbaselineskip}[0pt][0pt]{قنه} & \raisebox{-1\normalbaselineskip}[0pt][0pt]{ل} & \raisebox{-1\normalbaselineskip}[0pt][0pt]{كه} & \raisebox{-1\normalbaselineskip}[0pt][0pt]{\AbjadZero{}} & \raisebox{-1\normalbaselineskip}[0pt][0pt]{ش} & \raisebox{-1\normalbaselineskip}[0pt][0pt]{مظلم} \\
الكوكب الذي يتلوه على الضفيرة\textsuperscript{2} & قنط & م & كه & ل & ش & مظلم \\
\multicolumn{7}{|c|}{وتُسَمَّى هذه الثلثة ﴿الذَّوَائِبَ﴾} \\
\hline
\multicolumn{7}{|c|}{\large اسماء كواكب صورة العَذْراء والسُّنْبُلة} \\ \hline
الجَنوبيّ من الاثنَيْن اللذَيْن في رَأْس العَذْراء & قنز & ل & د & يه & ش & ه \\
الكوكب الشماليّ منه & قنح & ى & ه & ى & ش & ه \\
الذي على ظَهْرها عند الجَناح الأَيْسَر & قص & ى & ا & ى &  & ج \\
المقدَّم من الاربعة التي في الجَناح الايسر & قصا & كه & \AbjadZero{} & ى &  & د \\
الكوكب الذي يتلو هذا & قعب & ك & ب & ن &  & ج \\
الذي في الضِّلْع الأَيمن\textsuperscript{3} تحت الثَّدْي & قعب & ل & ح & ل &  & ج \\
المؤَخَّر من هذه الاربعة المذكورة & قعب & ى & ا & م &  & د \\
الذي على طَرَف القَدَم اليُسْرَى & را & ى & \AbjadZero{} & ل &  & د \\
الشماليّ الذي على طَرَف القَدَم اليُمْنَى & رج & ن & ط & ن &  & د \\
الذي على مَجْرَى ذَيْل السُّنْبُلة & قضز & ن & ز & ل &  & د \\
الذي في الجناح الشماليّ من الثلثة ﴿المُتَقَدِّم\newline للقَطاف﴾\textsuperscript{4} & \raisebox{-1\normalbaselineskip}[0pt][0pt]{رج} & \raisebox{-1\normalbaselineskip}[0pt][0pt]{ك} & \raisebox{-1\normalbaselineskip}[0pt][0pt]{يه} & \raisebox{-1\normalbaselineskip}[0pt][0pt]{ل} & \raisebox{-1\normalbaselineskip}[0pt][0pt]{ش} & \raisebox{-1\normalbaselineskip}[0pt][0pt]{ج} \\
النَّيِّر الذي على طَرَف يدها اليُسْرَى وهو السُّنْبُلة\newline ويُدْعَى ﴿السِّماك الأَعْزَل﴾ & \raisebox{-1\normalbaselineskip}[0pt][0pt]{رز} & \raisebox{-1\normalbaselineskip}[0pt][0pt]{ل} & \raisebox{-1\normalbaselineskip}[0pt][0pt]{ب} & \raisebox{-1\normalbaselineskip}[0pt][0pt]{\AbjadZero{}} & \raisebox{-1\normalbaselineskip}[0pt][0pt]{ش} & \raisebox{-1\normalbaselineskip}[0pt][0pt]{ا} \\
الكوكب الذي على منطقتها وعلى أَلْيَتها اليُمْنَى & رو & \AbjadZero{} & ح & م & ج & ج \\
\hline
\end{tabular}\end{Arabic}
\par\vspace{3mm}\noindent\rule{40mm}{.4pt}\par\vspace{1mm}{\small 1) Cod. \textarabic{المقبره} --- 2) Cod. \textarabic{المقبره} --- 3) Cod. \textarabic{الايسر} --- 4) Cod. \textarabic{للمطاف}}
\clearpage
{\large \textarabic{٢٦٠}}\par\vspace{2mm}
\begin{Arabic}\begin{tabular}{|p{84mm}|>{\centering\arraybackslash}p{9mm}|>{\centering\arraybackslash}p{9mm}|>{\centering\arraybackslash}p{9mm}|>{\centering\arraybackslash}p{9mm}|>{\centering\arraybackslash}p{14mm}|>{\centering\arraybackslash}p{17mm}|}\hline
\multirow{2}{*}{\parbox{80mm}{\centering من اسماء الكواكب الثابتة التي في صُوَر مِنْطَقَة البروج}} & \multicolumn{2}{c|}{الطول} & \multicolumn{2}{c|}{العرض} & \multirow{2}{*}{\parbox{13mm}{\centering\small علامات الجهة}} & \multirow{2}{*}{\parbox{16mm}{\centering\small مراتب العظمة}} \\ \cline{2-5}
 & {\small درج} & {\small دقائق} & {\small درج} & {\small دقائق} & & \\ \hline
\multicolumn{7}{|c|}{\large اسماء\textsuperscript{1} كواكب صورة المِيزَان} \\ \hline
الكوكب المُضِيء من كواكب الكَفَّة الجنوبيّة & رط & ى & ب & \AbjadZero{} & ش & ب \\
المُظْلِم الشماليّ من هذا الكوكب & رح & ى & ب & ل & ش & ه \\
النَّيِّر من الكواكب التي في طَرَف الكَفَّة الشماليّة & ريج & ه & ح & ك &  & ب \\
المُتَوَسِّط من التي في الكَفَّة الجنوبيّة & ريه & ى & ا & م &  & د \\
الكوكب التالي لهذا الكوكب & ريب & ل & ا & يه & ش & د \\
الكوكب المُتَوَسِّط من الكَفَّة الشماليّة & ريط & \AbjadZero{} & د & مه & ش & د \\
\hline
\multicolumn{7}{|c|}{\large اسماء كواكب صورة العَقْرَب} \\ \hline
الشماليّ من الثلثة التي بين عَيْنَيِ العَقْرَب & ركز & م & ا & ك & ش & ج \\
المتوسِّط من هذه الثلثة & ركو & ل & ا & ك & ج & ج \\
الكوكب الجنوبيّ من هذه الثلثة & ركو & ل & ه & \AbjadZero{} & ج & ج \\
المقدَّم من الثلثة المُضِيئَة التي في صَدْر العَقْرَب & رلا & ل & ج & مه &  & ج \\
﴿قَلْب العَقْرَب﴾ وهو الأَوْسَط منها الأَحمَر & رلج & ل & د & \AbjadZero{} &  & ب \\
الكوكب المُؤَخَّر من هذه الثلثة & رله & ل & ه & ل &  & ج \\
الكوكب الذي يتلو هذا في الخَرَزَةِ\textsuperscript{2} الأُولَى & رلط & م & يا & \AbjadZero{} &  & ج \\
الذي في الخَرَزَة الثانية & رم & \AbjadZero{} & يه & \AbjadZero{} &  & ج \\
المُضْعَف الشماليّ الذي في الخَرَزَة الثالثة & رما & ى & يح & م &  & د \\
الذي يتلوه في الخرزَة الرابعة & رمه & ك & يط & ل &  & د \\
التالي له في الخرزة الخامسة & رمط & ك & يح & ن & ج & ج \\
الذي يتلوه في الخرزة السادسة & رنا & م & يو & م & ج & ج \\
\hline
\end{tabular}\end{Arabic}
\par\vspace{3mm}\noindent\rule{40mm}{.4pt}\par\vspace{1mm}{\small 1) Inc. f. 232,r. --- 2) Cod. ubique \textarabic{الجزره}}
\clearpage
{\large \textarabic{٢٦١}}\par\vspace{2mm}
\begin{Arabic}\begin{tabular}{|p{84mm}|>{\centering\arraybackslash}p{9mm}|>{\centering\arraybackslash}p{9mm}|>{\centering\arraybackslash}p{9mm}|>{\centering\arraybackslash}p{9mm}|>{\centering\arraybackslash}p{14mm}|>{\centering\arraybackslash}p{17mm}|}\hline
\multirow{2}{*}{\parbox{80mm}{\centering من اسماء الكواكب الثابتة التي في صُوَر مِنْطَقَة البروج}} & \multicolumn{2}{c|}{الطول} & \multicolumn{2}{c|}{العرض} & \multirow{2}{*}{\parbox{13mm}{\centering\small علامات الجهة}} & \multirow{2}{*}{\parbox{16mm}{\centering\small مراتب العظمة}} \\ \cline{2-5}
 & {\small درج} & {\small دقائق} & {\small درج} & {\small دقائق} & & \\ \hline
الذي في الخرزة السابعة عند الشَّوْكَة وهي الشَّوْلة & رن & ى & يه & ى & ج & ج \\
المُؤَخَّر من الاثنيْن اللذيْن فِي الشوْلة & رمح & م & يج & ك &  & ج \\
المقدَّم منهما & رمح & ى & يج & ل & ج & ج \\
\multicolumn{7}{|c|}{وممَّا لَيس في صُورة} \\
وهو الشماليّ السَّحابيّ الذي يتلو الشَّوْلة & رنب & \AbjadZero{} & يج & ك & ج & غماميّ \\
\hline
\multicolumn{7}{|c|}{\large اسماء\textsuperscript{1} كواكب صورة الرَّامِي والقَوْس والسَّهْم مَعًا} \\ \hline
الكوكب الذي على زُجّ السَّهْم & رنه & م & و & ك & ج & ج\textsuperscript{2} \\
الذي على مِقْبَض يده اليُسْرَى & رنح & ل & و & ل & ج & ج \\
الكوكب الذي في الجَنوب من القَوْس & رنط & ى & ى & ن & ج & د \\
الجنوبيّ من الاثنَيْن اللذَيْنِ في الشِّقّ الشماليّ من\newline القَوْس & \raisebox{-1\normalbaselineskip}[0pt][0pt]{رص} & \raisebox{-1\normalbaselineskip}[0pt][0pt]{ى} & \raisebox{-1\normalbaselineskip}[0pt][0pt]{ا} & \raisebox{-1\normalbaselineskip}[0pt][0pt]{ل} & \raisebox{-1\normalbaselineskip}[0pt][0pt]{ج} & \raisebox{-1\normalbaselineskip}[0pt][0pt]{ج} \\
الشماليّ منهما وهو الذي في آخِر القَوْس & رنز & ن & ب & ن & ش & ج \\
الكوكب الذي على كَتِف الرامي & رصو & ل & ج & ى & ش & غماميّ \\
الذي بين يدَيْ هذا وهو في السَّهم & رصد & ى & د & مه & ش & د صغير \\
المُضْعَف السَّحابيّ الذي على عيْن الرامي & رصو & ك & \AbjadZero{} & مه & ش & ج \\
الذي على عُرْقُوب الرامي المقدَّم الأَيْسَر & رصح & ل & كج & \AbjadZero{} & ج & ج \\
المتوسِّط من الثلثة التي في ظَهْر الرامي وهو على\newline كَتِفه & \raisebox{-1\normalbaselineskip}[0pt][0pt]{رصح} & \raisebox{-1\normalbaselineskip}[0pt][0pt]{ل} & \raisebox{-1\normalbaselineskip}[0pt][0pt]{د} & \raisebox{-1\normalbaselineskip}[0pt][0pt]{ل} & \raisebox{-1\normalbaselineskip}[0pt][0pt]{ج} & \raisebox{-1\normalbaselineskip}[0pt][0pt]{ج} \\
التالي\textsuperscript{3} الذي تحت إبْطه ويَلِي الاوسط من التي في\newline ظَهْره & \raisebox{-1\normalbaselineskip}[0pt][0pt]{رصز} & \raisebox{-1\normalbaselineskip}[0pt][0pt]{ل} & \raisebox{-1\normalbaselineskip}[0pt][0pt]{و} & \raisebox{-1\normalbaselineskip}[0pt][0pt]{مه} & \raisebox{-1\normalbaselineskip}[0pt][0pt]{ج} & \raisebox{-1\normalbaselineskip}[0pt][0pt]{ج} \\
\hline
\end{tabular}\end{Arabic}
\par\vspace{3mm}\noindent\rule{40mm}{.4pt}\par\vspace{1mm}{\small 1) Inc. f. 232,v. --- 2) Cod. addit \textarabic{سحابى} --- 3) Cod. \textarabic{الثانى}}
\clearpage
{\large \textarabic{٢٦٢}}\par\vspace{2mm}
\begin{Arabic}\begin{tabular}{|p{84mm}|>{\centering\arraybackslash}p{9mm}|>{\centering\arraybackslash}p{9mm}|>{\centering\arraybackslash}p{9mm}|>{\centering\arraybackslash}p{9mm}|>{\centering\arraybackslash}p{14mm}|>{\centering\arraybackslash}p{17mm}|}\hline
\multirow{2}{*}{\parbox{80mm}{\centering من اسماء الكواكب الثابتة التي في صُوَر مِنْطَقَة البروج}} & \multicolumn{2}{c|}{الطول} & \multicolumn{2}{c|}{العرض} & \multirow{2}{*}{\parbox{13mm}{\centering\small علامات الجهة}} & \multirow{2}{*}{\parbox{16mm}{\centering\small مراتب العظمة}} \\ \cline{2-5}
 & {\small درج} & {\small دقائق} & {\small درج} & {\small دقائق} & & \\ \hline
الذي على رُكْبة الرامي من رِجْله اليُسْرَى & رصح & ى & يح & \AbjadZero{} & ج & ج \\
الذي على عُرْقوبه من رِجْله المقدَّمة & رنز & ل & يج & \AbjadZero{} & ج & ج \\
الذي في فَخِذه اليُسْرَى & رعح & ل & يج & ل &  & ج \\
الذي على ساق الرامي اليُمْنَى المُؤَخَّرة & رعز & ل & ك & ى &  & ج \\
الشماليّ من الاربعة التي في اصل أَلْيَته وهو\newline ﴿عُرْقوب الرامي﴾ & \raisebox{-1\normalbaselineskip}[0pt][0pt]{رعد} & \raisebox{-1\normalbaselineskip}[0pt][0pt]{\AbjadZero{}} & \raisebox{-1\normalbaselineskip}[0pt][0pt]{و} & \raisebox{-1\normalbaselineskip}[0pt][0pt]{ل} & \raisebox{-1\normalbaselineskip}[0pt][0pt]{ج} & \raisebox{-1\normalbaselineskip}[0pt][0pt]{ا} \\
الذي يتلوه في الخطّ الشماليّ & رف & ك & د & ن & ج & ه \\
\hline
\multicolumn{7}{|c|}{\large من اسْماء كواكب صورة الجَدْي} \\ \hline
المقدَّم من الثلثة التي في قَرْنه المُؤَخَّرة & رفح & ل & ز & ك & ش & ج \\
المتوسِّط من هذه الثلثة & رفح & ل & و & م & ش & و \\
الجنوبيّ من هذه الثلثة المذكورة & رفح & ل & ه & \AbjadZero{} &  & ج \\
الجنوبيّ من الثلثة التي في فَمِ الجَدْي & رض & ى & مه & مه &  & و \\
المقدَّم\textsuperscript{1} من الاثنَيْن الباقِيَيْن من الثلثة التي في فَم الجدي & رفط & ن & ا & مه &  & و \\
الكوكب الثالث الذي يتلوه في فَمه & رض & \AbjadZero{} & ا & ل & ش & و \\
الذي تحت رُكْبة الجَدْي اليُمْنَى & رضب & \AbjadZero{} & و & ل & ش & د \\
الذي على رُكْبته اليُسْرَى & رضب & ن & ح & م & ج & د \\
المقدَّم من الاثنيْن المتقارِبَيْن اللذَيْنِ في بَطْن الجدي & سا & ك & و & ن & ج & د \\
الكوكب التالي لهذا في بَطْنه & سا & ل & و & \AbjadZero{} &  & ه \\
المقدَّم من الاثنَيْن اللذيْن عند ذَنَب الجَدْي & سط & \AbjadZero{} & ب & ل & ج & ج \\
المقدَّم من الاربعة التي في شَمال ذَنَبه & سز & ل & ب & \AbjadZero{} & ج & ج \\
\hline
\end{tabular}\end{Arabic}
\par\vspace{3mm}\noindent\rule{40mm}{.4pt}\par\vspace{1mm}{\small 1) Inc. f. 233,r.}
\clearpage
{\large \textarabic{٢٦٣}}\par\vspace{2mm}
\begin{Arabic}\begin{tabular}{|p{84mm}|>{\centering\arraybackslash}p{9mm}|>{\centering\arraybackslash}p{9mm}|>{\centering\arraybackslash}p{9mm}|>{\centering\arraybackslash}p{9mm}|>{\centering\arraybackslash}p{14mm}|>{\centering\arraybackslash}p{17mm}|}\hline
\multirow{2}{*}{\parbox{80mm}{\centering من اسماء الكواكب الثابتة التي في صُوَر مِنْطَقَة البروج}} & \multicolumn{2}{c|}{الطول} & \multicolumn{2}{c|}{العرض} & \multirow{2}{*}{\parbox{13mm}{\centering\small علامات الجهة}} & \multirow{2}{*}{\parbox{16mm}{\centering\small مراتب العظمة}} \\ \cline{2-5}
 & {\small درج} & {\small دقائق} & {\small درج} & {\small دقائق} & & \\ \hline
المتوسِّط من هذه الثلثة الباقية & سط & ن & \AbjadZero{} & \AbjadZero{} & \AbjadZero{} & ه \\
الشماليّ من التي على طَرَف ذَنَبه & سح & ن & ب & ن & ش & ه \\
الكوكب الذي يتلوه & سط & ن & د & ك & ش & ه \\
المقدَّم من الذي في ظَهْر\textsuperscript{1} الجَدْي & سد & ل & د & مه & ج & د \\
المقدَّم من الكوكبَيْن اللذيْن في ظَهْره & رضز & ن & \AbjadZero{} & \AbjadZero{} & ج & د \\
الجَنوبيّ من الاربعة التي في شَمال ذَنَبه & سج & \AbjadZero{} & ج & \AbjadZero{} & ج & د \\
\hline
\multicolumn{7}{|c|}{\large من اسماء كواكب الدَّلْو والساقي} \\ \hline
الكوكب الذي على رَأْس الدَّلْو & سيا & ل & ه & مه & ش & ه \\
المُضِيء من الاثنَيْن اللذَيْنِ على كَتِف الساقي اليُمْنَى & سيز & ل & يا & \AbjadZero{} & ش & ج \\
الكوكب المُظْلِم الذي تحته & سيو & ك & ط & ن &  & ه \\
الذي في كَتِف الساقي اليُسْرَى & سز & ك & ح & ن &  & ج \\
الذي يتلو الثلثة التي في يده اليُسْرَى & رضح & ن & ه & ل &  & ج \\
المتوسِّط من هذه الثلثة & رضز & ك & ح & \AbjadZero{} &  & د \\
الكوكب المقدَّم من هذه الثلثة & رضه & ن & ح & ك &  & ج \\
الذي في فَخِذ الساق\textsuperscript{2} اليُمْنَى & سك & م & ح & ر &  & ج \\
الشمالي من الثلثة التي على طَرَف ذَنَبه\textsuperscript{3} & سكب & ن & ى & مه &  & ج \\
المقدَّم من الاثنيْن الجَنوبيَّيْن الباقِيَيْن & سكج & ى & ط & \AbjadZero{} &  & ج \\
الكوكب التالي لهذا & سكد & ل & ح & ل &  & ج \\
الجَنوبيّ من الاثنَيْن اللذيْن في ساقه اليُمْنَى & سكب & ن & ز & ل & ش & ج \\
المقدَّم\textsuperscript{4} من الخمسة التي في مَصَبّ الماء & سكو & ى & ب & \AbjadZero{} & ش & د \\
\hline
\end{tabular}\end{Arabic}
\par\vspace{3mm}\noindent\rule{40mm}{.4pt}\par\vspace{1mm}{\small 1) Error pro \textarabic{من الاثنين اللذين في صُلْب} --- 2) Error pro \textarabic{ذراع الساقي} --- 3) Error pro \textarabic{يده} --- 4) Inc. f. 233,v.}
\clearpage
{\large \textarabic{٢٦٤}}\par\vspace{2mm}
\begin{Arabic}\begin{tabular}{|p{84mm}|>{\centering\arraybackslash}p{9mm}|>{\centering\arraybackslash}p{9mm}|>{\centering\arraybackslash}p{9mm}|>{\centering\arraybackslash}p{9mm}|>{\centering\arraybackslash}p{14mm}|>{\centering\arraybackslash}p{17mm}|}\hline
\multirow{2}{*}{\parbox{80mm}{\centering من اسماء الكواكب الثابتة التي في صُوَر مِنْطَقَة البروج}} & \multicolumn{2}{c|}{الطول} & \multicolumn{2}{c|}{العرض} & \multirow{2}{*}{\parbox{13mm}{\centering\small علامات الجهة}} & \multirow{2}{*}{\parbox{16mm}{\centering\small مراتب العظمة}} \\ \cline{2-5}
 & {\small درج} & {\small دقائق} & {\small درج} & {\small دقائق} & & \\ \hline
الذي يتلوه في الجنوب & سكو & \AbjadZero{} & \AbjadZero{} & ى & ش & د \\
الذي يتلو هذا بعد المَقْبِض & سكح & ن & ا & ى & ج & د \\
الذي يتلوه ايضًا & سكا & م & ج & \AbjadZero{} & ج & د \\
الجَنوبيّ من التي في المَقْبِض & سلا & ل & ا & م &  & د \\
المؤَخَّر من التي في مَصَبّ الماء وهو ﴿فَم\textsuperscript{1} الحُوت\newline الجنوبيّ﴾ & \raisebox{-1\normalbaselineskip}[0pt][0pt]{سيح} & \raisebox{-1\normalbaselineskip}[0pt][0pt]{ى} & \raisebox{-1\normalbaselineskip}[0pt][0pt]{ك} & \raisebox{-1\normalbaselineskip}[0pt][0pt]{ك} & \raisebox{-1\normalbaselineskip}[0pt][0pt]{ج} & \raisebox{-1\normalbaselineskip}[0pt][0pt]{ا} \\
الشماليّ من الاثنَيْن اللذَيْن في الجَنوب\textsuperscript{2} من التي\newline في المَقْبِض & \raisebox{-1\normalbaselineskip}[0pt][0pt]{سل} & \raisebox{-1\normalbaselineskip}[0pt][0pt]{ى} & \raisebox{-1\normalbaselineskip}[0pt][0pt]{ج} & \raisebox{-1\normalbaselineskip}[0pt][0pt]{ل} & \raisebox{-1\normalbaselineskip}[0pt][0pt]{ج} & \raisebox{-1\normalbaselineskip}[0pt][0pt]{د} \\
المقدَّم من الاثنيْن المتقارِبَيْن اللذَيْن في ساق الساقي & سكب & ن & د & ل & ش & ج \\
الشماليّ من الاثنيْن اللذَيْن في جِهة الجنوب من\newline المقبض & \raisebox{-1\normalbaselineskip}[0pt][0pt]{سل} & \raisebox{-1\normalbaselineskip}[0pt][0pt]{ى} & \raisebox{-1\normalbaselineskip}[0pt][0pt]{ج} & \raisebox{-1\normalbaselineskip}[0pt][0pt]{ل} & \raisebox{-1\normalbaselineskip}[0pt][0pt]{ش} & \raisebox{-1\normalbaselineskip}[0pt][0pt]{ه} \\
الذي على أَلْيَة الساقي اليُمْنَى & سيط & ن & ج & ك & ج & ه \\
المتوسِّط من الثلثة التي في الانعطاف الأَوّل من\newline المقبض & \raisebox{-1\normalbaselineskip}[0pt][0pt]{سلج} & \raisebox{-1\normalbaselineskip}[0pt][0pt]{\AbjadZero{}} & \raisebox{-1\normalbaselineskip}[0pt][0pt]{يد} & \raisebox{-1\normalbaselineskip}[0pt][0pt]{مه} & \raisebox{-1\normalbaselineskip}[0pt][0pt]{ش} & \raisebox{-1\normalbaselineskip}[0pt][0pt]{د} \\
الاوّل من الثلثة التي في الانعطاف الثاني من المقبض & سكح & \AbjadZero{} & يد & مه & ج & د \\
الجنوبيّ من الاثنيْن اللذَيْن في أَلْيَته اليُسْرَى & سير & ن & يا & م & ج & د \\
المتوسِّط من الثلثة الأُخْرَى التي في الانعطاف الثاني\textsuperscript{3}\newline من المقبض & \raisebox{-1\normalbaselineskip}[0pt][0pt]{سكح} & \raisebox{-1\normalbaselineskip}[0pt][0pt]{م} & \raisebox{-1\normalbaselineskip}[0pt][0pt]{يه} & \raisebox{-1\normalbaselineskip}[0pt][0pt]{\AbjadZero{}} & \raisebox{-1\normalbaselineskip}[0pt][0pt]{ج} & \raisebox{-1\normalbaselineskip}[0pt][0pt]{د} \\
الشماليّ من الاثنيْن اللذيْن في ساقه اليُسْرَى تحت\newline الرُّكْبة & \raisebox{-1\normalbaselineskip}[0pt][0pt]{سيط} & \raisebox{-1\normalbaselineskip}[0pt][0pt]{\AbjadZero{}} & \raisebox{-1\normalbaselineskip}[0pt][0pt]{ط} & \raisebox{-1\normalbaselineskip}[0pt][0pt]{\AbjadZero{}} & \raisebox{-1\normalbaselineskip}[0pt][0pt]{ش} & \raisebox{-1\normalbaselineskip}[0pt][0pt]{ه} \\
\multicolumn{7}{|c|}{وممَّا لَيْسَ له في صُورة} \\
المقدَّم من الثلثة التي تتلو مَقْبِض الجَرَّة & سلز & ل & يه & ل & ج & ب كبير \\
\hline
\end{tabular}\end{Arabic}
\par\vspace{3mm}\noindent\rule{40mm}{.4pt}\par\vspace{1mm}{\small 1) Deest in cod. --- 2) Cod. \textarabic{الحوت} --- 3) Cod. \textarabic{الاول}}
\clearpage
{\large \textarabic{٢٦٥}}\par\vspace{2mm}
\begin{Arabic}\begin{tabular}{|p{84mm}|>{\centering\arraybackslash}p{9mm}|>{\centering\arraybackslash}p{9mm}|>{\centering\arraybackslash}p{9mm}|>{\centering\arraybackslash}p{9mm}|>{\centering\arraybackslash}p{14mm}|>{\centering\arraybackslash}p{17mm}|}\hline
\multirow{2}{*}{\parbox{80mm}{\centering من اسماء الكواكب الثابتة التي في صُوَر مِنْطَقَة البروج}} & \multicolumn{2}{c|}{الطول} & \multicolumn{2}{c|}{العرض} & \multirow{2}{*}{\parbox{13mm}{\centering\small علامات الجهة}} & \multirow{2}{*}{\parbox{16mm}{\centering\small مراتب العظمة}} \\ \cline{2-5}
 & {\small درج} & {\small دقائق} & {\small درج} & {\small دقائق} & & \\ \hline
الشماليّ من الاثنَيْن الباقِيَيْن & سم & ن & يد & م & ج & د كبير \\
الكوكب التالي لهما بأَثَر مقبض الجَرَّة & سم & ى & يح & يه & ج & د كبير \\
\hline
\multicolumn{7}{|c|}{\large من اسماء كواكب صُورتَيِ السَّمَكَتَيْن} \\ \hline
الذي في فَمِ الحُوت المقدَّم وهو الجنوبيّ & سلب & ن & ط & يه & ش & د كبير \\
الجنوبيّ من الاثنَيْن اللذيْن في رأْس هذا الحوت & سلد & ك & ز & ل & ش & ج \\
والشماليّ من هذَيْن الكوكبَيْن اللذيْن في رأْسه & سلز & ى & ط & ك &  & د \\
المقدَّم من الاثنيْن اللذيْن في ظَهْره & سلط & ك & ط & ل &  & د \\
الكوكب الثاني من الاثنيْن اللذيْن في ظهره & سما & ن & ك & ل &  & د \\
المقدَّم\textsuperscript{1} من الاثنيْن اللذيْن في بَطْن الحوت الجنوبيّ & سم & ى & ج & ل &  & د \\
الكوكب المؤَخَّر منهما & سمب & ن & ج & ل &  & د \\
الكوكب الذي في ذَنَب هذا الحوت & سمز & ى & و & ك &  & د \\
المقدَّم من التي في ذَنَبه في الخَيْط الكَتَّان & سنب & ى & ه & مه &  & و \\
الذي يتلو هذا الكوكب & سنه & ى & ج & مه &  & ه \\
المقدَّم من الثلثة النَّيِّرة التي بعد المتقدِّمة & سنح & ك & ب & يه &  & د \\
المتوسِّط منها & ا & م & ا & ى &  & د \\
الذي يتلو هذا من الثلثة & د & ى & و & \AbjadZero{} &  & د \\
الشماليّ من الاثنَيْن المتقارِبَيْن اللذَيْن في المَقْبِض & ج & ل & ب & \AbjadZero{} & ش & و \\
الكوكب الجنوبيّ منهما & د & ل & ه & \AbjadZero{} & ش & و \\
الاوسط من الثلثة التي بعد المَقْبِض & ط & ل & د & م & ج & د \\
الذي على مِرْبَط خَيْطَيِ الكَتَّان & يج & م & ح & ل & ش & ج \\
\hline
\end{tabular}\end{Arabic}
\par\vspace{3mm}\noindent\rule{40mm}{.4pt}\par\vspace{1mm}{\small 1) Inc. f. 234,r.}
\par\vfill\hfill 34
\clearpage
{\large \textarabic{٢٦٦}}\par\vspace{2mm}
\begin{Arabic}\begin{tabular}{|p{84mm}|>{\centering\arraybackslash}p{9mm}|>{\centering\arraybackslash}p{9mm}|>{\centering\arraybackslash}p{9mm}|>{\centering\arraybackslash}p{9mm}|>{\centering\arraybackslash}p{14mm}|>{\centering\arraybackslash}p{17mm}|}\hline
\multirow{2}{*}{\parbox{80mm}{\centering من اسماء الكواكب الثابتة التي في صُوَر مِنْطَقَة البروج}} & \multicolumn{2}{c|}{الطول} & \multicolumn{2}{c|}{العرض} & \multirow{2}{*}{\parbox{13mm}{\centering\small علامات الجهة}} & \multirow{2}{*}{\parbox{16mm}{\centering\small مراتب العظمة}} \\ \cline{2-5}
 & {\small درج} & {\small دقائق} & {\small درج} & {\small دقائق} & & \\ \hline
المقدَّم الشماليّ من مِرْبَط الكَتَّان & يا & م & ا & ك & ش & ه \\
المتوسِّط من الثلثة التي في المربط & يا & ن & ه & ك & ش & ج \\
الشماليّ من الاثنَيْن اللذيْن في فَم الحُوت الشماليّ & يج & ى & كا & مه &  & ه \\
الشماليّ من الثلثة التي على طَرَف الذَّنَب & يا & م & ط & \AbjadZero{} &  & د \\
المقدَّم من الثلثة التي على شَوْكة هذا الحوت & و & ن & يد & ك &  & د \\
المتوسِّط منها & ز & ن & يج & \AbjadZero{} &  & ج \\
المُؤَخَّر من هذه الثلثة & ح & ن & يب & \AbjadZero{} &  & د \\
الشماليّ من الاثنَيْن اللذيْن في بَطْنه & يج & ك & يز & \AbjadZero{} &  & و \\
الكوكب الجنوبيّ منهما & يا & \AbjadZero{} & يه & ك &  & د \\
الذي في شَوْكته المُؤَخَّرة عند ذَنَبه & يا & ى & يا & مه & ش & د \\
وممَّا ليس في صُوَرهما عند تحت المُرَبَّع & سنب & مو & ب & م & ج & د \\
\hline
\multicolumn{7}{|c|}{\parbox{168mm}{\centering\large من\textsuperscript{1} اسماء الثابتة التي في الصُّوَر الجنوبيّة عن مِنْطَقَة البروج لسنة \AbjadOver{اقضا}\textsuperscript{2} لذي القَرْنَيْن}} \\ \hline
\multicolumn{7}{|c|}{\large من كواكب قيطس وهو سَبُع البَحْر} \\ \hline
الذي على طَرَف أَنْف سَبُع البَحْر وهو قيطس\textsuperscript{3} & كح & ن & ز & مه & ج & ج \\
المُؤَخَّر من الثلثة التي في حُلْقومه على طَرَف لَحيه & كح & ن & يب & ك & ج & ج \\
المتوسِّط منها وهو في وسط فَمه & كج & ن & يا & ل &  & ج \\
المقدَّم من الثلثة التي على ذَقْنه & كا & م & يد & \AbjadZero{} &  & ج \\
الذي على جَبِينه فَوْقَ عَيْنه & كا & ك & ح & مه & ج & د \\
الجنوبيّ من الشِّقّ المُؤَخَّر منه & يح & ى & كز & ل & ج & ج \\
\hline
\end{tabular}\end{Arabic}
\par\vspace{3mm}\noindent\rule{40mm}{.4pt}\par\vspace{1mm}{\small 1) Inc. f. 234,v. --- 2) Apud Maghrebinos 1191. --- 3) Cod. \textarabic{قيطوس}}
\clearpage
{\large \textarabic{٢٦٧}}\par\vspace{2mm}
\begin{Arabic}\begin{tabular}{|p{84mm}|>{\centering\arraybackslash}p{9mm}|>{\centering\arraybackslash}p{9mm}|>{\centering\arraybackslash}p{9mm}|>{\centering\arraybackslash}p{9mm}|>{\centering\arraybackslash}p{14mm}|>{\centering\arraybackslash}p{17mm}|}\hline
\multirow{2}{*}{\parbox{80mm}{\centering من اسماء الكواكب الثابتة التي في الصُّوَر الجنوبيّة عن مِنْطَقَة البروج}} & \multicolumn{2}{c|}{الطول} & \multicolumn{2}{c|}{العرض} & \multirow{2}{*}{\parbox{13mm}{\centering\small علامات الجهة}} & \multirow{2}{*}{\parbox{16mm}{\centering\small مراتب العظمة}} \\ \cline{2-5}
 & {\small درج} & {\small دقائق} & {\small درج} & {\small دقائق} & & \\ \hline
المتوسِّط من الثلثة التي في جَسَده\textsuperscript{1} & ج & ى & كه & ك & ج & ج \\
الشماليّ من هذه الثلثة & و & ى & ك & \AbjadZero{} & ج & ج \\
المؤَخَّر من الكوكبَيْن اللذيْن في ذَنَبه & \AbjadZero{} & ن & يه & ك &  & ج \\
المقدَّم منهما & سمو & ى & يه & ل &  & ج \\
الشماليّ الذي في الشِّقّ المقدَّم منه & سن & ل & يج & \AbjadZero{} &  & ه كبير \\
الشماليّ الذي عند الاثنَيْن اللذيْن في طَرَف ذَنَبه & سمه & ن & ط & م & ج & ج كبير \\
الكوكب الجنوبيّ الذي على طَرَف ذَنَبه & سمه & ن & ك & م & ج & د صغير \\
\hline
\multicolumn{7}{|c|}{\large ومن كواكب الجَبَّار} \\ \hline
السَّحابيّ الذي في رَأْسه & صح & ى & يو & ل & ج & غماميّ \\
النَّيِّر الذي على كَتِفه اليُمْنَى وهو ﴿مَنْكِب الجوزاء﴾ & عج & ى & يز & \AbjadZero{} & ج & ا صغير \\
الذي على كَتِفه اليُسْرَى & صه & ى & يز & ل &  & ب كبير \\
والذي\textsuperscript{2} تحت هذا الذي في كَتِفه اليُسْرَى & صو & ك & ح & \AbjadZero{} &  & د صغير \\
الشماليّ من التِّسْعة التي في الجِلْد الذي بيده اليُسْرَى & صا & م & ح & \AbjadZero{} &  & د \\
الكوكب الذي على مِرْفَقه الأَيْمن & عه & ل & يد & ل &  & د \\
السادس منها الذي في الشمال & نو & \AbjadZero{} & يه & ن &  & ج \\
السابع الذي بعد هذا في الشمال & نو & \AbjadZero{} & يز & ى &  & ج \\
الثامن الذي بعده في الشمال & نو & ل & كا & ك &  & ج \\
التاسع وهو في الجنوب من التي في الجلد وتسمّى\newline ﴿المِجْلَد﴾\textsuperscript{3} & \raisebox{-1\normalbaselineskip}[0pt][0pt]{نز} & \raisebox{-1\normalbaselineskip}[0pt][0pt]{ل} & \raisebox{-1\normalbaselineskip}[0pt][0pt]{كط} & \raisebox{-1\normalbaselineskip}[0pt][0pt]{ل} & \raisebox{-1\normalbaselineskip}[0pt][0pt]{ج} & \raisebox{-1\normalbaselineskip}[0pt][0pt]{ج} \\
المقدَّم من الثلثة التي في المِنْطقة & صو & ل & كد & ى & ج & ب \\
\hline
\end{tabular}\end{Arabic}
\par\vspace{3mm}\noindent\rule{40mm}{.4pt}\par\vspace{1mm}{\small 1) Cod. \textarabic{جيده} --- 2) Haec linea, in textu omissa, in margine addita legitur. --- 3) Ex coniectura; cod. \textarabic{المحله}}
\clearpage
{\large \textarabic{٢٦٨}}\par\vspace{2mm}
\begin{Arabic}\begin{tabular}{|p{84mm}|>{\centering\arraybackslash}p{9mm}|>{\centering\arraybackslash}p{9mm}|>{\centering\arraybackslash}p{9mm}|>{\centering\arraybackslash}p{9mm}|>{\centering\arraybackslash}p{14mm}|>{\centering\arraybackslash}p{17mm}|}\hline
\multirow{2}{*}{\parbox{80mm}{\centering من اسماء الكواكب الثابتة التي في الصُّوَر الجنوبيّة عن مِنْطَقَة البروج}} & \multicolumn{2}{c|}{الطول} & \multicolumn{2}{c|}{العرض} & \multirow{2}{*}{\parbox{13mm}{\centering\small علامات الجهة}} & \multirow{2}{*}{\parbox{16mm}{\centering\small مراتب العظمة}} \\ \cline{2-5}
 & {\small درج} & {\small دقائق} & {\small درج} & {\small دقائق} & & \\ \hline
الاوسط منها & صح & ل & كد & ل & ج & ب \\
المُؤَخَّر من الثلثة التِي في المنطقة & صط & ك & كه & م & ج & ب \\
الذي عند مِقْبَض سَيْفه & صه & \AbjadZero{} & كه & ن &  & ج \\
الشماليّ\textsuperscript{1} من الثلثة المُضافة التي عند رأْس السيْف & صز & ن & ك & م &  & د \\
المتوسِّط من هذه الثلثة & صز & ن & كط & ى &  & ج صغير \\
المقدَّم من هذه الثلثة وهو الجنوبيّ منها & صز & ن & كط & ك &  & ج \\
النَّيِّر الذي على طَرَف رِجْله اليُسْرَى وهو ﴿رِجْل\newline الجَوْزاء﴾ & \raisebox{-1\normalbaselineskip}[0pt][0pt]{صا} & \raisebox{-1\normalbaselineskip}[0pt][0pt]{\AbjadZero{}} & \raisebox{-1\normalbaselineskip}[0pt][0pt]{لا} & \raisebox{-1\normalbaselineskip}[0pt][0pt]{ل} &  & \raisebox{-1\normalbaselineskip}[0pt][0pt]{ا} \\
الكوكب الشماليّ من هذا وهو فوْق عُرْقوبه & صب & ى & ل & يه &  & د \\
الخارج الذي تحت عَقِبه اليُسْرَى\textsuperscript{2} & صد & ل & لا & ى & ج & د \\
الذي تحت رُكْبته اليُمْنَى المُؤَخَّرة & عا & ك & لج & ل & ج & ج كبير \\
\hline
\multicolumn{7}{|c|}{\large ومن كواكب النَّهْر} \\ \hline
الذي على طَرَف رِجْل الجَبَّار وهو على رأْس النَّهْر & صط & ل & ل & ن & ج & د كبير \\
النَّيِّر الكبير وهو آخِر كواكب النهر & ىا & ك & لج & ل & ج & ا \\
الذي يتلو الاربعة التي في القَطْع\textsuperscript{3} & لح & ى & لب & ن &  & ج \\
الثالث بين يَدَيِ الاوسط & له & ك & كح & ن &  & ج \\
المقدَّم من الاربعة & لج & ى & كح & \AbjadZero{} &  & ج \\
المؤَخَّر من الاربعة & كح & ل & كه & ل &  & ج \\
الكوكب الثالث الذي قبل هذا الرابع & كج & ك & كج & ل & ج & ج \\
الذي في مُنْعَطَف النهر وهو على آخِر صَدْر قيطس\textsuperscript{4} & ىو & ك & لب & ى & ج & د \\
\hline
\end{tabular}\end{Arabic}
\par\vspace{3mm}\noindent\rule{40mm}{.4pt}\par\vspace{1mm}{\small 1) Inc. f. 235,r. --- 2) Cod. \textarabic{الايسر} --- 3) Cod. \textarabic{الخط} --- 4) Cod. \textarabic{صبر القلوس}}
\clearpage
{\large \textarabic{٢٦٩}}\par\vspace{2mm}
\begin{Arabic}\begin{tabular}{|p{84mm}|>{\centering\arraybackslash}p{9mm}|>{\centering\arraybackslash}p{9mm}|>{\centering\arraybackslash}p{9mm}|>{\centering\arraybackslash}p{9mm}|>{\centering\arraybackslash}p{14mm}|>{\centering\arraybackslash}p{17mm}|}\hline
\multirow{2}{*}{\parbox{80mm}{\centering من اسماء الكواكب الثابتة التي في الصُّوَر الجنوبيّة عن مِنْطَقَة البروج}} & \multicolumn{2}{c|}{الطول} & \multicolumn{2}{c|}{العرض} & \multirow{2}{*}{\parbox{13mm}{\centering\small علامات الجهة}} & \multirow{2}{*}{\parbox{16mm}{\centering\small مراتب العظمة}} \\ \cline{2-5}
 & {\small درج} & {\small دقائق} & {\small درج} & {\small دقائق} & & \\ \hline
الشماليّ الذي في الشِّقّ المقدَّم من الاربعة التي في\newline الجَوَاز\textsuperscript{1} & \raisebox{-1\normalbaselineskip}[0pt][0pt]{لب} & \raisebox{-1\normalbaselineskip}[0pt][0pt]{ل} & \raisebox{-1\normalbaselineskip}[0pt][0pt]{ما} & \raisebox{-1\normalbaselineskip}[0pt][0pt]{ك} & \raisebox{-1\normalbaselineskip}[0pt][0pt]{ج} & \raisebox{-1\normalbaselineskip}[0pt][0pt]{د} \\
المُؤَخَّر من هذه الاربعة & له & ن & مج & ك & ج & د \\
المؤَخَّر من الاثنَيْن اللذيْن بعد العَطْف\textsuperscript{2} ويُسَمَّى\newline ﴿التربه﴾ & \raisebox{-1\normalbaselineskip}[0pt][0pt]{لط} & \raisebox{-1\normalbaselineskip}[0pt][0pt]{ك} & \raisebox{-1\normalbaselineskip}[0pt][0pt]{نج} & \raisebox{-1\normalbaselineskip}[0pt][0pt]{ن} &  & \raisebox{-1\normalbaselineskip}[0pt][0pt]{د} \\
المؤَخَّر من الثلثة التي بعد الاثنَيْن & كط & \AbjadZero{} & نج & لا\textsuperscript{3} &  & د \\
الشَّماليّ من الاثنَيْن المتقابِلَيْن & مه & ك & صج & ك & ج & د \\
الجَنوبيّ منهما & مو & ى & نا & ن & ج & د \\
\hline
\multicolumn{7}{|c|}{\large ومن كواكب الأَرْنَب} \\ \hline
الشماليّ من الاربعة التي في ظَهْرها في الشِّقّ المقدَّم & ص & ل & له & \AbjadZero{} & ج & ه \\
الشماليّ من الشقّ المؤَخر منها & صب & ل & له & م & ج & ه \\
الكوكب\textsuperscript{4} الذي في ذَقَن\textsuperscript{5} الأَرْنَب & صا & \AbjadZero{} & لو & ل &  & ه \\
الذي بين يَدَيْ طَرَف الرِّجْل المقدَّمة اليُسْرَى منها & ص & ك & لط & يه &  & د كبير \\
الذي على وَسَط جَسَدها & نز & ك & مه & يه &  & د كبير \\
الكوكب الذي تحت بَطْنها & صا & \AbjadZero{} & ما & ل &  & ج \\
الشماليّ من الاثنَيْن اللذيْن في الرِّجْلَيْن المُؤَخَّرَتَيْن & صه & ل & مد & ن &  & ج \\
الكوكب الجنوبيّ منهما & عب & ى & مد & \AbjadZero{} &  & د كبير \\
الكوكب الذي على ظَهْرها & ع & ى & مه & مه & ج & د كبير \\
الكوكب الذي على طَرَف أَلْيَتها & عا & ى & لح & م & ج & د كبير \\
\hline
\end{tabular}\end{Arabic}
\par\vspace{3mm}\noindent\rule{40mm}{.4pt}\par\vspace{1mm}{\small 1) Ex coniectura; cod. \textarabic{الجوزا} --- 2) Ex coniectura; cod. \textarabic{الصحرا} --- 3) Scilicet \textarabic{\AbjadZero{}} --- 4) Incipit f. 235,v. --- 5) Cod. \textarabic{ذنب}}
\clearpage
{\large \textarabic{٢٧٠}}\par\vspace{2mm}
\begin{Arabic}\begin{tabular}{|p{84mm}|>{\centering\arraybackslash}p{9mm}|>{\centering\arraybackslash}p{9mm}|>{\centering\arraybackslash}p{9mm}|>{\centering\arraybackslash}p{9mm}|>{\centering\arraybackslash}p{14mm}|>{\centering\arraybackslash}p{17mm}|}\hline
\multirow{2}{*}{\parbox{80mm}{\centering من اسماء الكواكب الثابتة التي في الصُّوَر الجنوبيّة عن مِنْطَقَة البروج}} & \multicolumn{2}{c|}{الطول} & \multicolumn{2}{c|}{العرض} & \multirow{2}{*}{\parbox{13mm}{\centering\small علامات الجهة}} & \multirow{2}{*}{\parbox{16mm}{\centering\small مراتب العظمة}} \\ \cline{2-5}
 & {\small درج} & {\small دقائق} & {\small درج} & {\small دقائق} & & \\ \hline
\multicolumn{7}{|c|}{\large اسماء كواكب صورة الكَلْب} \\ \hline
الكوكب الذي على أذْنَيْه & عج & ن & لح & ى & ج & د كبير \\
الكوكب الذي في رأْسه & ض & ل & له & \AbjadZero{} & ج & د \\
النَّيِّر الذي في فَم الكَلْب وهو ﴿الشِّعْرَى اليمانيّة﴾ & ضب & ل & لو & ل &  & ه \\
الذي على طَرَف رِجْله المقدَّمة & قح & ن & نط & ى &  & ا \\
الذي في أَصْل فَخِذه اليُسْرَى & قب & ى & ما & ك &  & ج \\
الذي في اصل فخذه اليمنى & ضز & ن & مح & مه &  & ج صغير \\
الذي تحت بَطْنه فيما بين فَخِذَيْه & ضد & ن & نا & ل &  & ج \\
الذي على رِجْله اليُمْنَى & ق & ن & نج & مه &  & ج \\
الكوكب الذي في ذَنَبه & قج & ك & ل & م &  & غماميّ صغير \\
\multicolumn{7}{|c|}{وممَّا لَيْسَ\textsuperscript{1} له في صورة} \\
المؤَخَّر من الاثنَيْن المُضِيئَيْن & ع & ى & نط & م & ج & ب \\
المقدَّم منهما & صز & ى & نز & م & ج & ب \\
\hline
\multicolumn{7}{|c|}{\large من كواكب مُقَدَّم الكَلْب} \\ \hline
الذي في مقدَّم الكَلْب وهو في عُنُقه & ضو & ى & ند & \AbjadZero{} & ج & د \\
المُضِيء التالي الذي في حَقه وهو ﴿الشِّعْرَى الشَّامِيَّة﴾ & ق & ك & نو & ى & ج & ا \\
\hline
\multicolumn{7}{|c|}{\parbox{168mm}{\centering\large بَقيّة\textsuperscript{2} كواكب صورة اقراطير\textsuperscript{3} وهو الكَأْس}} \\ \hline
الكوكب الذي على تدوير فَم الكَأْس الشَّمالي & قص & ل & يج & م & ج & د \\
\hline
\end{tabular}\end{Arabic}
\par\vspace{3mm}\noindent\rule{40mm}{.4pt}\par\vspace{1mm}{\small 1) Cod. \textarabic{ليست} --- 2) Inc. f. 236,r. Folium igitur in codice desideratur, quo continebantur stellae Navis, Hydrae, et primae stellae Crateris. --- 3) Cod. \textarabic{ابروطورش}}
\clearpage
{\large \textarabic{٢٧١}}\par\vspace{2mm}
\begin{Arabic}\begin{tabular}{|p{84mm}|>{\centering\arraybackslash}p{9mm}|>{\centering\arraybackslash}p{9mm}|>{\centering\arraybackslash}p{9mm}|>{\centering\arraybackslash}p{9mm}|>{\centering\arraybackslash}p{14mm}|>{\centering\arraybackslash}p{17mm}|}\hline
\multirow{2}{*}{\parbox{80mm}{\centering من اسماء الكواكب الثابتة التي في الصُّوَر الجنوبيّة عن مِنْطَقَة البروج}} & \multicolumn{2}{c|}{الطول} & \multicolumn{2}{c|}{العرض} & \multirow{2}{*}{\parbox{13mm}{\centering\small علامات الجهة}} & \multirow{2}{*}{\parbox{16mm}{\centering\small مراتب العظمة}} \\ \cline{2-5}
 & {\small درج} & {\small دقائق} & {\small درج} & {\small دقائق} & & \\ \hline
الكوكب الذي على أُذْن الكأس الجَنوبيّة & قعا & ك & يو & ى & ج & د \\
الذي على أُذْن الكأس الشماليّة & قصب & ن & يا & ل & ج & ه \\
\hline
\multicolumn{7}{|c|}{\large اسماء كواكب صورة الغُرَاب} \\ \hline
الذي في مِنقار الغُراب عند الشُّجاع & قفو & ل & كا & م & ج & ج \\
الذي في عُنُقه قريب من رَأْسه & قفح & ل & يط & م & ج & ج \\
الذي في الجَناح المقدَّم الأَيْمَن & قعد & م & يد & ن &  & ج \\
المقدَّم من الاثنَيْن اللذين في الجَناح المُؤَخَّر & قعز & ل & يب & ل &  & ج \\
الكوكب المُؤَخَّر منهما & قعح & ى & يا & مه & ج & د \\
الذي على طَرَف رِجْله عند الشُّجاع & قفا & م & يح & ن & ج & د \\
\hline
\multicolumn{7}{|c|}{\large من اسماء كواكب قنطاورس\textsuperscript{1} وهو صورة انسان وفَرَس ويُسَمَّى الظُّلْمان} \\ \hline
الجَنوبيّ من الاربعة التي في رَأْس قنطاورس\textsuperscript{2} & را & ك & كا & م & ج & ه كبير \\
الكوكب الشماليّ منها\textsuperscript{3} & را & ى & يح & ن & ج & ه كبير \\
المقدَّم من الاثنَيْن الأَوْسَطَيْن الباقِيَيْن & ر\AbjadZero{} & ك & ك & ل &  & د كبير \\
المُؤَخَّر من هذَيْن الاثنَيْن وهو الثاني من الاربعة & را & ى & ك & \AbjadZero{} &  & ه كبير \\
الذي على كَتِفه المقدَّمة اليُسْرَى & قضز & ك & كه & م &  & ج \\
الذي على كَتِفه المُؤَخَّرة اليُمْنَى & رو & ل & كب & ل &  & ج \\
الذي على رَأْس الغِراس\textsuperscript{4} من الاثنَيْن الباقِيَيْن & ريج & ى & يح & يه &  & د كبير \\
المقدَّم من الثلثة التي على الشِّقّ الأَيْمَن منه & رد & ل & كح & ك & ج & د كبير \\
الاوسط منها & ره & ى & كط & ك & ج & د كبير \\
\hline
\end{tabular}\end{Arabic}
\par\vspace{3mm}\noindent\rule{40mm}{.4pt}\par\vspace{1mm}{\small 1) Cod. \textarabic{فنطارس} --- 2) Cod. \textarabic{فيطارس} --- 3) Cod. \textarabic{منهما} --- 4) Cod. \textarabic{الفارس}}
\clearpage
{\large \textarabic{٢٧٢}}\par\vspace{2mm}
\begin{Arabic}\begin{tabular}{|p{84mm}|>{\centering\arraybackslash}p{9mm}|>{\centering\arraybackslash}p{9mm}|>{\centering\arraybackslash}p{9mm}|>{\centering\arraybackslash}p{9mm}|>{\centering\arraybackslash}p{14mm}|>{\centering\arraybackslash}p{17mm}|}\hline
\multirow{2}{*}{\parbox{80mm}{\centering من اسماء الكواكب الثابتة التي في الصُّوَر الجنوبيّة عن مِنْطَقَة البروج}} & \multicolumn{2}{c|}{الطول} & \multicolumn{2}{c|}{العرض} & \multirow{2}{*}{\parbox{13mm}{\centering\small علامات الجهة}} & \multirow{2}{*}{\parbox{16mm}{\centering\small مراتب العظمة}} \\ \cline{2-5}
 & {\small درج} & {\small دقائق} & {\small درج} & {\small دقائق} & & \\ \hline
المؤَخَّر من هذه الثلثة & رو & ك & كح & \AbjadZero{} & ج & د كبير \\
الذي على ذِراعه اليُمْنَى & رز & ل & كو & ل & ج & كبير \\
الذي على ساعِده الايمن & ريد & لا\textsuperscript{1} & كه & يه &  & ج \\
المُضِيء الذي في أَصْل جَنْبه الايسر & رط & ى & لج & ل &  & ج \\
المؤَخَّر من الكواكب المُظْلِمة الشماليَّة منه & رح & ن & لا & \AbjadZero{} &  & د \\
الكوكب\textsuperscript{2} المقدَّم من الكواكب المُظْلِمة الشماليّة منه & رح & \AbjadZero{} & ل & ك &  & ه \\
المؤَخَّر من الثلثة التي على خاصِرته على يمينه & قضج & \AbjadZero{} & م & \AbjadZero{} &  & ج \\
المقدَّم من الاثنَيْن المتقارِبَيْن اللذَيْن في فَخِذ الفَرَس اليُمْنَى & قضج & ل & مو & ى &  & ج \\
المقدَّم من الاثنَيْن اللذَيْن تحت بَطْن الفَرَس & رز & ك & مج & \AbjadZero{} &  & ب \\
الكوكب المؤَخَّر منهما & رح & ن & مج & مه &  & ج \\
الذي في فَخِذه اليُمْنَى قريب من الرِّجْل & را & ى & نا & ى &  & ب \\
الذي على قَدَمه اليُسْرَى على الحافِر & رب & ك & نه & ك &  & ب \\
النَّيِّر الذي على طَرَف رِجْله اليُمْنَى المقدَّمة وهو\newline ﴿رِجْل الفرس﴾\textsuperscript{3} & \raisebox{-1\normalbaselineskip}[0pt][0pt]{قضط} & \raisebox{-1\normalbaselineskip}[0pt][0pt]{ل} & \raisebox{-1\normalbaselineskip}[0pt][0pt]{ما} & \raisebox{-1\normalbaselineskip}[0pt][0pt]{ى} &  & \raisebox{-1\normalbaselineskip}[0pt][0pt]{ا} \\
الكوكب الذي على رُكْبته اليُسْرَى من الرِّجْل & ريو & \AbjadZero{} & مه & ك &  & ب \\
الذي على عُرْقوبه الايمن & رز & ل & نا & م &  & ب \\
السادس من التي على الرِّجْل المُؤَخَّرة اليُمْنَى & ره & ن & مط & ى & ج & د \\
الذي تحت وَسَط فَخِذه اليُسْرَى & قضز & ل & نه & ى & ج & د \\
\hline
\multicolumn{7}{|c|}{\large اسماء كواكب صورة السَّبُع} \\ \hline
الكوكب الذي على طَرَف رِجْل السَّبُع المُؤَخَّرة & رىط & ى & كد & ن & ج & ج \\
\hline
\end{tabular}\end{Arabic}
\par\vspace{3mm}\noindent\rule{40mm}{.4pt}\par\vspace{1mm}{\small 1) I. e. \textarabic{\AbjadZero{}} --- 2) Inc. f. 236,v. --- 3) Forte \textarabic{الفرس} error pro \textarabic{قنطاورس}; cf. p. \textarabic{٢٧٥}, ult. l.}
\clearpage
{\large \textarabic{٢٧٣}}\par\vspace{2mm}
\begin{Arabic}\begin{tabular}{|p{84mm}|>{\centering\arraybackslash}p{9mm}|>{\centering\arraybackslash}p{9mm}|>{\centering\arraybackslash}p{9mm}|>{\centering\arraybackslash}p{9mm}|>{\centering\arraybackslash}p{14mm}|>{\centering\arraybackslash}p{17mm}|}\hline
\multirow{2}{*}{\parbox{80mm}{\centering من اسماء الكواكب الثابتة التي في الصُّوَر الجنوبيّة عن مِنْطَقَة البروج}} & \multicolumn{2}{c|}{الطول} & \multicolumn{2}{c|}{العرض} & \multirow{2}{*}{\parbox{13mm}{\centering\small علامات الجهة}} & \multirow{2}{*}{\parbox{16mm}{\centering\small مراتب العظمة}} \\ \cline{2-5}
 & {\small درج} & {\small دقائق} & {\small درج} & {\small دقائق} & & \\ \hline
الذي على وَسَط فَخِذه المؤَخَّرة & ريو & م & كط & ى & ج & ج \\
الشماليّ من الثلثة التي على طَرَف ذَنَبه & ريد & ى & كط & ك & ج & د كبير \\
الجَنوبيّ من الاثنَيْن اللذَيْن في رِجْله المقدَّمة & ريح & ك & يا & ل &  & د كبير \\
الشماليّ من الاثنَيْن اللذَيْن في عُنُقه & رل & ل & يه & ك & ج & د كبير \\
الشماليّ من الاثنَيْن اللذَيْن في رِجْله المقدَّمة & ريز & م & ى & \AbjadZero{} & ج & د كبير \\
\hline
\multicolumn{7}{|c|}{\large من اسماء كواكب صورة المِجْمَرَة وتسمَّى المُرِيح} \\ \hline
الشماليّ من الاثنَيْن اللذَيْن في أَسْفَل المِجْمَرَة & رمح & ن & كب & م & ج & ه \\
الذي في وَسَط رَأْسها وهو موضِع النار & رمز & ك & كو & ل & ج & د كبير \\
الجَنوبيّ من الاثنين المتقارِبَيْن اللذَيْن في لَهَب النار & رمو & ك & لد & ى &  & د كبير \\
الذي على طَرَف اللَّهَب من لِسان النار & رمب & \AbjadZero{} & لد & \AbjadZero{} & ج & د \\
 & رند & ى & كا & مه & ج & د \\
\hline
\multicolumn{7}{|c|}{\large اسماء\textsuperscript{1} كواكب صورة الإكلِيل الجَنوبيّ} \\ \hline
المقدَّم من الستّة التي في تَقويس الإكلِيل الجَنوبيّ & رص & ك & كا & ل & ج & د \\
الكوكب الرابع من هذه الستّة & رصو & \AbjadZero{} & ك & \AbjadZero{} & ج & د \\
الذي يتلوه وهو بين يَدَيْ رُكْبَة الرامي & رصز & ك & يح & ل &  & ه \\
المُضِيء الذي يتلو هذا من الشمال & رصح & ى & يز & ى &  & ب \\
الشماليّ من هذا المُضِيء & رعب & \AbjadZero{} & يو & \AbjadZero{} &  & د \\
المقدَّم من الاثنين المُظْلِمَيْن & رصه & ن & يد & ن & ج & و \\
الكوكب الباقي من المُظْلِمَيْن & رصج & \AbjadZero{} & يد & م & ج & ه \\
\hline
\end{tabular}\end{Arabic}
\par\vspace{3mm}\noindent\rule{40mm}{.4pt}\par\vspace{1mm}{\small 1) Inc. f. 237,r.}
\par\vfill\hfill 35
\clearpage
{\large \textarabic{٢٧٤}}\par\vspace{2mm}
\begin{Arabic}\begin{tabular}{|p{84mm}|>{\centering\arraybackslash}p{9mm}|>{\centering\arraybackslash}p{9mm}|>{\centering\arraybackslash}p{9mm}|>{\centering\arraybackslash}p{9mm}|>{\centering\arraybackslash}p{14mm}|>{\centering\arraybackslash}p{17mm}|}\hline
\multirow{2}{*}{\parbox{80mm}{\centering من اسماء الكواكب الثابتة التي في الصُّوَر الجنوبيّة عن مِنْطَقَة البروج}} & \multicolumn{2}{c|}{الطول} & \multicolumn{2}{c|}{العرض} & \multirow{2}{*}{\parbox{13mm}{\centering\small علامات الجهة}} & \multirow{2}{*}{\parbox{16mm}{\centering\small مراتب العظمة}} \\ \cline{2-5}
 & {\small درج} & {\small دقائق} & {\small درج} & {\small دقائق} & & \\ \hline
المقدَّم من هذا ايضًا & رص & ل & يه & ن & ج & ه \\
الكوكب الباقي الجنوبيّ من هذا & رص & ك & يح & ن & ج & ه \\
\hline
\multicolumn{7}{|c|}{\large اسماء كواكب صورة الحُوت الجَنوبيّ} \\ \hline
الكوكب الذي في فَمِ الحُوت الجنوبيّ على طَرَف الماء & سيا & ن & ك & ك & ج & د \\
التالي لهذا الكوكب & سيه & ك & كب & يه & ج & د \\
الثالث المؤَخَّر من هذه الثلثة المقدَّمة & سيو & ل & كب & ى &  & د \\
الكوكب الذي في حُلْقوم الحوت & سيه & ل & يو & ى &  & د كبير \\
الجنوبيّ الذي في الشَّوْكة الجنوبيّة & سو & ك & ط & ل &  & ه \\
المؤَخَّر من الاثنَيْن اللذَيْن في بَطْنه & سيب & ك & يه & ى &  & د \\
المقدَّم منهما & سى & \AbjadZero{} & يد & م &  & ج \\
المؤَخَّر من الثلثة التي في الشَّوْكة الشماليّة & سو & ك & يب & \AbjadZero{} &  & د \\
المتوسِّط من هذه الثلثة & سج & \AbjadZero{} & يو & ل &  & د \\
المقدَّم من هذه الثلثة & سب & ى & يح & ى &  & د \\
الذي على طَرَف ذَنَبه & سا & ك & كب & يه &  & د \\
\multicolumn{7}{|c|}{وعنده مِمَّا لَيْسَ\textsuperscript{1} له في صُورة} \\
المقدَّم من الثلثة المُضِيئة & رفط & ى & كب & ك &  & ج صغير \\
الاوسط من هذه الثلثة & رصب & ى & كب & ى &  & ج صغير \\
المؤَخَّر من هذه الثلثة & رضه & ى & كا & \AbjadZero{} & ج & ج صغير \\
الكوكب المُظْلِم الذي بين يَدَيْه & رضج & ى & ك & ن & ج & د \\
\hline
\end{tabular}\end{Arabic}
\par\vspace{3mm}\noindent\rule{40mm}{.4pt}\par\vspace{1mm}{\small 1) Cod. \textarabic{ليست}}
\end{document}
